\documentclass{report}

\input{~/dev/latex/template/preamble.tex}
\input{~/dev/latex/template/macros.tex}

\title{\Huge{}}
\author{\huge{Nathan Warner}}
\date{\huge{}}
\fancyhf{}
\rhead{}
\fancyhead[R]{\itshape Warner} % Left header: Section name
\fancyhead[L]{\itshape\leftmark}  % Right header: Page number
\cfoot{\thepage}
\renewcommand{\headrulewidth}{0pt} % Optional: Removes the header line
%\pagestyle{fancy}
%\fancyhf{}
%\lhead{Warner \thepage}
%\rhead{}
% \lhead{\leftmark}
%\cfoot{\thepage}
%\setborder
% \usepackage[default]{sourcecodepro}
% \usepackage[T1]{fontenc}

% Change the title
\hypersetup{
    pdftitle={DC Circuits}
}

\begin{document}
    % \maketitle
        \begin{titlepage}
       \begin{center}
           \vspace*{1cm}
    
           \textbf{DC Circuits}
    
           \vspace{0.5cm}
            
                
           \vspace{1.5cm}
    
           \textbf{Nathan Warner}
    
           \vfill
                
                
           \vspace{0.8cm}
         
           \includegraphics[width=0.4\textwidth]{~/niu/seal.png}
                
           Computer Science \\
           Northern Illinois University\\
           United States\\
           
                
       \end{center}
    \end{titlepage}
    \tableofcontents
    \pagebreak 
    \unsect{Basic quantities}
    \begin{itemize}
        \item \textbf{Planetary model of the atom}: One of the major issues with the planetary model is the idea that electrons whirl around the nucleus in stable, planet-like orbits. That's simply not true. First, the electron inhabits a region of 3D space and does not simply move through a plane       
            \bigbreak \noindent 
            Second, due to the Heisenberg Uncertainty Principle, we can't precisely plot the position and trajectory of a given electron. The best we can do is make a plot of where the electron is likely to be. This is called a probability contour
            \bigbreak \noindent 
            To understand this concept, imagine that you could record the position of a specific electron relative to the nucleus at a specific time. A moment later you record its new position, a moment after that you record the next position, and so on for thousands of measurements. If you attempted to plot them all, you would wind up with a cloud of dots around the nucleus. This cloud is referred to as an orbital. You wouldn't know how the electron got from one position to the next but you would get a general idea of where it was likely to be. Ultimately, an orbital looks nothing like a planetary orbit.
        \item \textbf{Shells, subshells, and orbitals}: The electron cloud is referred to as an \textbf{orbital.}
            \bigbreak \noindent 
            We note that there are several potential orbitals. Due to quantum physics, only certain orbitals are allowed. The permissible electron energy levels are first grouped into shells, then subshells and finally orbitals. It is important to remember that orbitals indicate the electron energy level. That is, a higher orbital implies a higher energy level. Further, orbitals fill in first from lowest energy level to highest energy level.
        \item \textbf{Charge and current}: Charge is an attractive force denoted by the letter $Q$ and has units of coulombs.
            \bigbreak \noindent 
            Electrons are negatively charged and protons are positively charged. All electrons and protons exhibit the same magnitude of charge
            \bigbreak \noindent 
            Further, opposite charges attract while like charges repel, similar to the poles of a magnet.
        \item \textbf{Moving charge}: It is possible to move charge from one point to another. The rate of charge movement over time is called current. It is denoted by the letter I and has units of amps
            \begin{align*}
                1 \text{ amp} = 1 \text{ coulomb} / 1 \text{ second}
            .\end{align*}
            Thus, current for current $I$, charge $Q$, and time $t$,
            \begin{align*}
                I = \frac{Q}{t}
            .\end{align*}
        \item \textbf{Energy and Voltage}:  Energy is defined as the ability to do work. It is denoted by the letter W. The basic unit is the joule
            \bigbreak \noindent 
            If we were to move a charge from one point to another (for example, separating an electron from an atom), we would have to expend energy to do so
            \bigbreak \noindent 
            \fig{.6}{./figures/1.png}
            We would say that B has a higher electric potential than A. In other words, there is a potential difference between B and A. We refer to this change as voltage, denoted $V$, with units volts. One volt is one joule per coulomb.
            \begin{align*}
                V = \frac{W}{Q} 
            .\end{align*}
            \begin{itemize}
                \item V is the voltage in volts,
                \item W is the energy in joules,
                \item Q is the charge in coulombs
            \end{itemize}
            As you might guess, the bigger the charge to be moved, the greater the energy required
        \item \textbf{Power and efficiency}: Power is the rate of energy usage
            \begin{align*}
                P = \frac{W}{t}
            .\end{align*}
            Measured in watts.
            \begin{align*}
                1 \text{watt} = 1 \text{ joule } / 1 \text{ second}
            .\end{align*}
        \item \textbf{The power law}: Recall that
            \begin{align*}
                I = \frac{Q}{t}, \quad V = \frac{W}{Q}, \quad \text{and} \quad P = \frac{W}{t}
            .\end{align*}
            So,
            \begin{align*}
                P = \frac{VQ}{t} = IV
            .\end{align*}
            We call $P = VI$ the \textbf{power law}.
        \item \textbf{Efficiency}: Efficiency is the ratio of useful output power to applied power expressed as a percentage
            \begin{align*}
                \eta = \frac{P_{\text{out}}}{P_{\text{in}}} \times 100\%
            .\end{align*}
            If a device draws 200 watts of power and has a useful output of 120 watts, the efficiency is
            \begin{align*}
                \eta = \frac{120}{200} \cdot 100\% = 60\%
            .\end{align*}
            In this case, the device is wasting 40\% of the input power, or 80 watts.
        \item \textbf{amp-hour}: An amp hour (Ah, or ampere-hour) is a unit of measurement used to describe the total electric charge or capacity of a battery.
            \bigbreak \noindent 
            One amp hour means a battery can deliver a steady current of 1 amp for 1 hour.
            \bigbreak \noindent 
            The lifespan of a batter is roughly the amp-hour rating over the current draw.
            \begin{align*}
                \text{Lifespan} \approx \frac{Ah}{I}
            .\end{align*}
            Amp-hour is calculated by multiplying current by time
            \begin{align*}
                Ah = It
            .\end{align*}
            It measures the total electric charge a battery can hold or deliver.
        \item \textbf{Resistance and conductance}:  Resistance is a measure of how difficult it is to establish a flow of current
            \bigbreak \noindent 
            Sometimes it is more convenient to use the inverse view of this phenomenon, and instead of referring to how difficult it is to establish a current, we would be interested in how easy it is to establish a current. This is called conductance and it is the reciprocal of resistance. Conductance is denoted by the letter G and has units of siemens (S).
            \begin{align*}
                R = \frac{1}{G}, \quad G = \frac{1}{R}
            .\end{align*}
        \item \textbf{Resistivity}: The measure of a material's inherent and general ability to restrict current flow is called resistivity. Resistivity is denoted by the Greek letter $\rho$ 
            \bigbreak \noindent 
            All other factors being equal, the higher the resistivity, the greater will be the resistance. Further, the greater the length of the material that the current must pass through, the greater the resistance. Finally, as the surface area of the cross section (i.e, the face) grows, the resistance decreases.
            \begin{align*}
                R = \frac{\rho \cdot l}{A}
            .\end{align*}
            \begin{itemize}
                \item $R$ is the resistance in ohms,
                \item $\rho$ is the resistivity,
                \item $l$ is the length of the material,
                \item $A$ is area of the face ($h$ times $w$).
            \end{itemize}
            Resistivity is often specified in ohm-centimeters with the length and area similarly specified in centimeters and square centimeters
        \item \textbf{Resistors}: Resistors are \textbf{linear bilateral devices}. Being linear, their current-voltage relation can be drawn as a straight line. Bilateral means that the polarity of voltage or direction of current will not matter. In other words, unlike a battery, these devices cannot be inserted into a circuit backwards because either orientation works the same way. If the horizontal axis is voltage and the vertical is current, then the slope of the line yields the conductance
            \bigbreak \noindent 
            If the x-axis is labeled voltage, and y-axis is labeled current, the steeper the line, the lower its resistance
            \bigbreak \noindent 
            General purpose resistors use a color code to denote their value and tolerance. Typically, this will involve four color stripes: two for the precision/mantissa, one for the power of ten, and the fourth to indicate the tolerance. For higher precision, a five stripe version may be used with the first three denoting the precision/ mantissa. Alternately, high precision resistors may have their nominal value printed directly on them
            \bigbreak \noindent 
            \fig{.6}{./figures/2.png}
            \bigbreak \noindent 
            \fig{.6}{./figures/3.png}
            \bigbreak \noindent 
            Note that the fourth band is spaced away from the others to avoid accidentally reversing the order.
        \item \textbf{Overvoltage (overload voltage)}: Overvoltage is a condition where the electrical supply voltage rises above the safe design limit or rated nominal value of a circuit or device
        \item \textbf{Force Sensing Resistor (FSR)}: The force sensing resistor consists of two layers of material with a nominal separation distance. It is presented as a flat membrane that is usually round or rectangular, and perhaps with an adhesive backing. 
            \bigbreak \noindent 
            With no force applied, the device shows an extremely high resistance, well into the megohms. As force is applied to the surface, the two layers come into better contact which decreases the net resistance
        \item \textbf{Photoresistor (light dependent resistor (LDR))}:  As their name implies, photoresistors are sensitive to changes in light level. They are also called LDRs, short for Light Dependent Resistor
            \bigbreak \noindent 
             In total darkness the device exhibits a very high resistance. As light levels increase, the resistance decreases.
        \item \textbf{Thermistor}: A thermistor is a device whose resistance is a function of temperature. These devices are available in two basic types. Either PTC, for Positive Temperature Coefficient; or NTC, for Negative Temperature Coefficient. PTC devices show an increase in resistance as temperature increases and NTC devices show a decrease in resistance as temperatures rise. Ideally, these are linear relationships with the plots showing straight lines.
        \item \textbf{Varistor}: Varistors are used as limiting devices, primarily to suppress unwanted voltage spikes in electronic equipment. The varistor is a unique device in that it has a highly nonlinear current-voltage characteristic. 
            \bigbreak \noindent 
            Remember, when plotted with the voltage on the horizontal axis, the slope of the line represents the conductance. Consequently, the varistor shows a region of near zero conductance or extremely high resistance (the horizontal section), and two sections that are nearly vertical, which indicate extremely high conductance or near zero resistance. This characteristic allows the varistor to act as a limiting device.
            \bigbreak \noindent 
            \fig{.6}{./figures/4.png}
        \item \textbf{Schematic symbols}: Clearly, it would not be practical to create circuit drawings (schematics) using pictures of the actual devices. Instead, we use simple schematic symbols to represent them. There are two widely used schemes: ANSI (American National Standards Institue) and IEC (International Electrotechnical Commission). For many electronic components the ANSI and IEC versions are the same, but there are notable differences. ANSI tends to dominate in North America (surprise!) while IEC tends to dominate elsewhere. 
            \bigbreak \noindent 
            Figure 2.38 shows the symbols for DC current and voltage sources. The long bar at the top of the voltage source is the positive terminal while the shorter bar at the bottom denotes the negative terminal. For the current source, the arrow points in the direction of conventional current flow. If the source is adjustable, it will often be shown with a diagonal arrow drawn through it. In some schematics, a DC voltage source will simply be drawn as a “connection dot” or node with the voltage labeled. It is assumed that the other end of the source is tied back to the system common or ground.
            \bigbreak \noindent 
            \fig{.6}{./figures/5.png}
            \bigbreak \noindent 
            Speaking of ground, there are three different symbols used for a circuit's common connection point. These are: earth ground, chassis ground and signal or digital ground. These are shown in Figure 2.39. Earth ground is used when the circuit common is tied back to true earth (e.g., to the third pin on an AC receptacle). Chassis ground is a more generic term which would represent a common reference point that did not go back to true earth (e.g., a portable device). Finally, digital or signal ground is used to distinguish between common points where they might be separated for noise or interference concerns (e.g., a sensitive analog signal common separate from a digital logic common).
            \bigbreak \noindent 
            \fig{.6}{./figures/6.png}
            \bigbreak \noindent 
            The symbols for resistors and other resistive devices vary considerably between ANSI and IEC versions. ANSI uses a zig-zag line while IEC uses a simple rectangle
            \bigbreak \noindent 
            One final note about component values: As decimal points are easy to lose on photocopies or small display screens, engineering notation multipliers are sometimes placed in the position of a decimal point. For example, a 4.7 k$\Omega$ resistor may be listed as simply 4k7.

    \end{itemize}

    \pagebreak 
    \unsect{Series Resistive Circuits}: 
    \begin{itemize}
        \item \textbf{Intro}: Benjamin Franklin (pictured in Figure 3.1) began experimenting with the phenomenon of electricity in 1746. In 1752 he performed his famous kite experiment proving that lightning is a form of electricity by capturing charge from storm clouds in a leyden jar (an early form of an electrical charge storage device). 
            \bigbreak \noindent 
            At this time the modern concept of an atomic model with electrons and protons did not exist and electricity was conceived of as a sort of fluid. Franklin surmised that the “electrical flow” moved from positive to negative. This idea was accepted and became the conventional view. Today we call this idea \textbf{conventional current flow}. In this model, current flows from a more positive voltage to a less positive voltage. We know now that the electron is the charge carrier in metals and the electrons travel in the reverse direction. Essentially, Franklin guessed wrong. Electrons move from a lower potential to a higher potential. We call this model \textbf{electron flow}. For most work, engineers and technicians use conventional flow, although in some cases, such as the explanation of semiconductors, electron flow is easier to visualize for some people. In short, conventional flow exists for historical reasons, and it is the model used for most analyses, including this text.
            \bigbreak \noindent 
            You might think the current direction would make a big difference in an analysis; after all, it certainly makes a big difference if you drive a car in the wrong direction. It turns out that both forms will achieve the desired results, we just have to be consistent with the usage.
            \bigbreak \noindent 
            To better understand this, consider that the movement of a net negative charge in one direction can be thought of as a movement of a net positive charge in the other direction. That is, the movement of an electron creates a “hole” where it used to be and that hole is net positive.
            \bigbreak \noindent 
            \fig{.6}{./figures/7.png}
            \bigbreak \noindent 
            Here we start at the top with a tube of identical marbles, all pushed to the right. In each step below it we move a marble to the left, mimicking the flow of electrons in a circuit. When we reach the bottom, each marble has been pushed left by one place. We can also arrive at the bottom drawing by simply taking the right-most marble in the top drawing and inserting it to the extreme left by jumping over the other three marbles. Here's the important bit: instead of imagining the marbles moving left, we can also think in terms of “negative space” and imagine the empty slot moving to the right. That's hole flow. The two views are functionally identical as they lead to the same outcome.
        \item \textbf{Circuits}: The word circuit comes from the Latin root circ, meaning “ring” or “around”. An electrical circuit consists of at least one ring or loop through which current flows. For example, if we have a battery attached to a lamp the current exits the battery, flows through the lamp, and then returns to the other side of the battery creating a loop or completed circuit. Without a path back to the battery, current will not flow. Thus, if we cut one of the wires connecting the battery and lamp, there is no path for current and no current flows. This is referred to as an open circuit and is a common fault that occurs when electronic systems are dropped or struck forcefully. Obviously, this will tend to render the circuit unusable. In this example, the lamp will not illuminate.
            \bigbreak \noindent 
            The opposite of the open circuit is the short circuit. In a short circuit, an unintended alternate path for current flow exists and this also can create a malfunction. In the case of our battery and lamp, a short circuit can occur if a piece of wire or metal accidentally fell across the terminals of the lamp. The current would then have a high conductance (i.e., low resistance) path around the lamp. The vast majority of the current would take this low resistance path instead of the higher resistance path presented by the lamp. The result would be that the lamp would not light. In either the open or short case, the light does not function but there is an important difference: for the short circuit, excessive current will flow out of the battery because there is little to resist the flow of current. Thus, the battery will be drained very quickly. For the open, no current flows and thus the battery is not drained.
        \item \textbf{Series connections}: In general, a series connection is any connection of components configured such that the current through each component must be the same
            \bigbreak \noindent 
            Inside each of the lettered boxes would be a component such as a resistor or a voltage source. Note that for each component, there is one entrance point and one exit point. No matter which box you pick, the current flowing through it must be the same as the current flowing into the next box or out of the preceding box, and it doesn't matter if you follow this path in a clockwise or counterclockwise fashion. Thus, this entire configuration is a series connection. The idea that current is consistent throughout should be selfevident. After all, the only way the current through, say, item B could be different from that flowing through item C or D is if some of it somehow “disappeared” along the way by magic. It is important to remember that consistent current is the hallmark feature defining a series connection:
            \bigbreak \noindent 
            It is possible that only a portion of a circuit exhibits a series connection.
            \bigbreak \noindent 
            Consider the more complex diagram presented in Figure 3.5. Some of these items are in series and some are not. For example, items A and B are in series with each other but not in series with the remaining items. Why? Because if we imagine the current flowing through A and then through B they must be the same, however, once beyond B, the current could split and flow down other branches: a portion entering C, a portion entering D and the remainder flowing into E. On the other hand, items C and F are in series with each other because whatever current is flowing through one of them must be flowing through the other. Thus, items A and B are in series with each other and items C and F are in series with each other, although all four are not in series as a group.
            \bigbreak \noindent 
            \fig{.6}{./figures/9.png}
            \bigbreak \noindent 
            It is possible that no two items in a circuit are strictly in series. Further, just because the currents through two items happen to be the same, that does not necessarily mean they are in series. Identical currents could be just a by-product of the component values chosen
            \bigbreak \noindent 
            Even if items D and E in Figure 3.5 have the same numeric value for current, we would not say that they are in series, anymore than we would say that any two people with the same last name would have to be siblings.
        \item \textbf{Combining Series Components}: Typically, a series connection will include multiple resistors. 
            \bigbreak \noindent 
            Recall that resistors add
            \begin{align*}
                R = \frac{\rho \ell}{ A}
            \end{align*}
            resistance. If we consider two identical resistors placed in series, one after the other the effective length would double while keeping the resistivity and area unchanged. The combined result would be a doubling of the resistance of just one of them. 
            \bigbreak \noindent 
            \fig{.6}{./figures/10.png}
            \bigbreak \noindent 
             If we then generalize this to two arbitrary resistors of identical resistivity and area, the lengths would dictate the resistance of each, and the combined lengths would then reflect the resistance of the pair.
             \begin{align*}
                 R &= \frac{\rho(\ell_{1} + \cdots + \ell_{n})}{A} = \frac{\rho}{A}(\ell_{1} + \cdots + \ell_{n}) \\
             .\end{align*}
             Thus we find that the equivalent resistance of a group of series resistors is their sum:
             \begin{align*}
                 R_{\text{total}} = R_{1} + R_{2} + \cdots + R_{n}
             .\end{align*}
             Consequently, as resistors in series add, total resistance may be found by summing the individual resistors.
             \bigbreak \noindent 
             \fig{1}{./figures/11.png}
             \bigbreak \noindent 
             Multiple voltage sources in series may also be added, however, polarities must be considered as opposing voltages partially cancel each other (i.e., adding a negative).
             \bigbreak \noindent 
             \fig{.6}{./figures/12.png}
             \bigbreak \noindent 
             If we use point b as our reference, by inspection the top of the 12 volt source is 12 volts above point b (reminder, the long bar denotes the positive terminal). Also, by inspection, the right side of the 3 volt source (point a) is negative with respect to its left side. As the left side of this source is connected to the positive terminal of the 12 volt source, then it too must be 12 volts above point b. As its right side is 3 volts less than this side, point a must be 3 volts less than 12 volts, or 9 volts above point b. Thus 
             \begin{align*}
                 V_{ab} = 9 volts.
             .\end{align*}
             Chasing this further, if the 3 volt source had been flipped so that it had the same polarity as the 12 volt source (positive toward the right) then $V_{ab} = 15$ volts. If the 12 volt source had been flipped (positive toward bottom) with the 3 volt source as drawn originally, then $V_{ab} = -15$ volts. If both sources had been flipped then $V_{ab} = -9$ volts. Finally, if point a had been taken as the reference then these four potentials would have the opposite polarity because, by definition, $V_{ba} = -V_{ba}$
             \bigbreak \noindent 
             In contrast to voltage sources, differing current sources are not placed in series as they would each attempt to establish a different series current, a practical impossibility
        \item \textbf{Generality}: 
            \begin{align*}
                \text{Effect} = \frac{\text{cause}}{\text{opposition}}
            .\end{align*}
        \item \textbf{Ohm's law}: We have
            \begin{align*}
                V = IR
            .\end{align*}
            Or,
            \begin{align*}
                R = \frac{V}{I}
            .\end{align*}
            Note that from this, we get
            \begin{align*}
                P = I^{2}R = \frac{V^{2}}{R}
            .\end{align*}
        \item \textbf{Kirchoff's Voltage Law (KVL)}: Along with Ohm's law, the key law governing series circuits is Kirchhoff's voltage law, or KVL. Named after nineteenth century German physicist Gustav Kirchhoff, this law states that the sum of voltage rises and voltage drops around a series loop must equal zero (the rises and drops having opposite polarities). Alternately, it may be reworded as the sum of voltage rises around a series loop must equal the sum of voltage drops
            \begin{align*}
                \sum V \uparrow\; = \sum V \downarrow
            .\end{align*}
        \item \textbf{Voltage divider rule (VDR)}: An outgrowth of KVL is the voltage divider rule (VDR). In a series connection, the current is the same through each component. Thus, the voltage drops in a series connection must be directly proportional to the size of the resistances: the larger the resistor, the larger its voltage, and the larger its share of the total voltage applied to the series connection. Thus, the voltage across any resistor must equal the net supplied voltage times the ratio of the resistor of interest to the total resistance:
            \begin{align*}
                V_{R_{x}} = E \cdot \frac{R_{x}}{R_{\text{total}}}
            .\end{align*}
            This is just two Ohm's law calculations:
            \begin{align*}
                V_{\text{total}} = E = IR_{\text{total}} \implies I = \frac{E}{R_{\text{total}}}
            .\end{align*}
            Then, for a particular resistor $R_{x}$:
            \begin{align*}
                V_{R_{x}} = IR_{x} = \frac{E}{R_{\text{total}}} \cdot R_{x} = E \cdot \frac{R_{x}}{R_{\text{total}}}
            .\end{align*}
            It is worth pointing out that “the resistor of interest” can, in fact, be the sum of multiple resistors in series.
            \bigbreak \noindent 
            While VDR is not required for any particular analysis, it serves two purposes: first, it saves some time because it skips over the computation of current, and second, it reinforces the ideal of a proportional division of voltage in a series connection. For example, if there are two resistors in series and one of them is twice the size of the other, then it must be the case that the larger resistor sees twice the voltage of the smaller resistor.
        \item \textbf{Series analysis}: Ohm's law, Kirchhoff's voltage law, the voltage divider rule and series component combinations are the tools we will use to solve general series circuit problems. There are multiple techniques for analyzing these circuits:
            \begin{itemize}
                \item If all voltage source and resistor values are given, the circulating current can be found by dividing the equivalent voltage by the sum of the resistances. Once the current is found, Ohm's law can be used to find the voltage drops across individual resistors. Alternately, the individual voltage drops can be found using the voltage divider rule. At that point, power law can be used to find the power dissipation in each resistor or the power developed by the source(s).
                \item If the circuit uses a current source instead of a voltage source, then the circulating current is known and the voltage drop across any resistor may be determined directly using Ohm's law.
                \item If the problem concerns determining resistance values, the basic idea will be to use these rules in reverse. For example, if a resistor value is needed to set a specific current, the total required resistance can be determined from this current and the given voltage supply. The values of the other series resistors can then be subtracted from the total, yielding the required resistor value. Similarly, if the voltages across two resistors are known, as long as one of the resistance values is known, the other resistance can be determined using either the voltage divider rule or Ohm's law. 
            \end{itemize}
        \item \textbf{Node voltage}: A node voltage is the voltage at a given point with respect to some other point, normally ground
            \bigbreak \noindent 
            Note that as we treat connecting wires ideally (having no resistance), a node is the entire connection wire, not a specific location on it. Normally, simulators will automatically number each node in a circuit with ground being node 0. This method skips the entire business of inserting virtual meters in the circuit and reduces schematic clutter.
        \item \textbf{Voltage rise and voltage drop}: A voltage rise occurs when you move from a lower electric potential to a higher electric potential:
            \bigbreak \noindent 
            A voltage drop occurs when you move from a higher electric potential to a lower electric potential:
        \item \textbf{The Importance of Polarity}: A key item of importance when analyzing series circuits, or indeed any electrical circuit, is noting the polarity of the voltages. 
            \bigbreak \noindent 
            Determining polarity for voltage sources is easy as their polarity is fixed (positive is always at the “long bar” and negative at the “short bar”). The polarity of any resistor will depend on the direction of current flow. To avoid confusion we will use the following standard:
            \begin{itemize}
                \item Where conventional current enters a resistor, we label this point with a plus sign, and where the current leaves, a minus sign.
                \item During analysis, moving from plus to minus is a voltage drop. That is, we are moving from a more positive or higher potential to a less positive or lower potential, thus “dropping” voltage.
                \item Similarly, traversing from minus to plus is a voltage rise.
            \end{itemize}
            Remember, Kirchhoff's voltage law states that the sum of voltage rises around a series loop must equal the sum of voltage drops.
            \bigbreak \noindent 
            Along with determining the voltage across single resistors, we are often interested in determining the voltage between arbitrary points in a circuit. These may span several components. It is imperative that we know the polarities of the individual voltages in order to determine the voltage between any two circuit points.
            \bigbreak \noindent 
            \fig{.6}{./figures/13.png}
            \bigbreak \noindent 
            Here, the two voltage sources, $E_{1}$ and $E_{2}$, aid each other because they are both trying to establish a clockwise current. This produces a net voltage of $E_{1}+E_{2}$. Their polarities are fixed. The current direction will be as drawn because conventional current flows out of the positive terminal of a voltage source. The value of this current will be 
            \begin{align*}
                I = \frac{E_{1} + E_{2}}{R_{1} + R_{2}}
            \end{align*}
            via Ohm's law.
            \bigbreak \noindent 
            Knowing the direction of current, the polarities of the voltage drops across the two resistors may be found. The point where the current enters the resistor is positive, and negative where it exits. Once these are labeled, either Ohm's law or the voltage divider rule can be used to determine the voltages developed across the two resistors. Note that KVL states that the sum of these two resistor voltages must equal the net supplied voltage, or $E_{1}+E_{2}$ in this case.
            \bigbreak \noindent 
            To determine the voltage from point $a$ to point $c$, or $V_{ac}$, we start at point $a$ and sum the voltage drops and rises until we get to point $c$. Here, the voltage across $R_{1}$ shows up as a drop
            \bigbreak \noindent 
            The result is the voltage across $R_{1}$ minus the voltage across $E_{2}$.
            \bigbreak \noindent 
            Depending on the specific values, the sum may wind up being either a positive or negative value. If it's positive, this signifies that point $a$ is at a higher potential than is point $c$. Conversely, if it's negative, this signifies that point $a$ is at a lower potential than is point $c$.
            \bigbreak \noindent 
            To illustrate the importance of voltage source polarity, consider the series circuit shown in Figure 3.22.
            \bigbreak \noindent 
            \fig{.6}{./figures/14.png}
            \bigbreak \noindent 
            The two voltage sources oppose each other because their positive terminals are connected through \(R_1\), rather than one source’s positive terminal aiding the other source.
            \bigbreak \noindent 
            Assume $E_{2} > E_{1}$. 
            \bigbreak \noindent 
            Following the indicated current direction around the loop:
            \begin{itemize}
                \item Passing through \(E_2\) from \(c\) to \(b\) produces a voltage rise of \(+E_2\).
                \item Passing through \(E_1\) from \(a\) downward produces a voltage drop of \(-E_1\).
            \end{itemize}
            Therefore, the net voltage driving the circuit is
            \begin{align*}
                E_{\text{net}} = E_{2} - E_{1}
            .\end{align*}
            The current is therefore
            \begin{align*}
                I = \frac{E_{2} - E_{1}}{R_{1} + R_{1}}
            .\end{align*}
            Both resistors are in series, so the same current flows through them.
            \bigbreak \noindent 
            Because \(E_2\) is the stronger source, conventional current leaves its positive terminal at \(b\). It then travels:
            \begin{align*}
               b \to R_{1} \to a \to E_{1} \to R_{2} \to c \to E_{2} 
            .\end{align*}
            Thus, current passes:
            \begin{itemize}
                \item Through \(R_1\) from right to left.
                \item Through \(E_1\) from its positive terminal to its negative terminal.
                \item Through \(R_2\) from left to right.
            \end{itemize}
            A source delivers energy when conventional current leaves its positive terminal. Current leaves the positive terminal of \(E_2\), so \(E_2\) supplies power.
            \bigbreak \noindent 
            Current enters the positive terminal of $E_{1}$, so $E_{1}$ absorbs power.
            \begin{align*}
                P_{E_{1}} = E_{1}I
            .\end{align*}
            If \(E_1\) represents a rechargeable battery, that absorbed power corresponds to charging. If it is an ideal voltage source, “charging” simply means that it is absorbing energy.
            \begin{itemize}
                \item \(E_{1} < E_{2}\): current flows as drawn.
                \item \(E_1>E_2\): current flows in the opposite direction.
                \item \(E_1=E_2\): the ideal circuit has zero current.
            \end{itemize}
        \item \textbf{Potentiometers and Rheostats}: Some applications require the use of variable resistances. One example is a basic volume control, but other examples include equipment calibration and user adjustable parameters such as the bass, treble and balance controls on a home audio system. Adjustable resistances fall under two categories, potentiometers and rheostats. Generally, these would be the same physical device, with the application and usage determining its name. 
            \bigbreak \noindent 
            A potentiometer is a resistor with three leads. Along with the typical end leads, there is a third lead referred to as the wiper or tap. This third lead is connected to a small metallic arm that makes contact with the resistive material and which can be moved along the length of the material. Potentiometers, or pots for short, often are formed in a circular shape (rotary travel) but they can also be in the form of a straight line (linear travel)
            \bigbreak \noindent 
            The larger circular units with shafts and the straight unit at the top are designed to be adjusted by the end user (appropriately stylish knobs being attached to said shafts, of course), while the smaller units are intended to be adjusted only by a technician (note the small adjustment screw slots intended to receive fine screwdrivers). The smaller units are designed to be mounted directly on a printed circuit board while the user adjustable units are mounted on a front panel, and have wiring lugs instead of straight leads.
            \bigbreak \noindent 
            A pot can be visualized as two resistors in sequence with a tap between them. The total resistance never changes. What does change is the ratio of the two pieces.
            \begin{align*}
                R_{AB} = R_{AW} + R_{WB}
            .\end{align*}
            \bigbreak \noindent 
            \fig{.6}{./figures/15.png}
            \bigbreak \noindent 
        \item \textbf{Ganged potentiometers}: As a result, potentiometers are very useful as controllable voltage dividers (hence the name). Sometimes it is desirable to control two signals simultaneously from a single knob (e.g., one volume knob for a hi-fi stereo controls both the left and right channels). For this application, multiple resistive elements can be stacked and use a common shaft to control all of the wiper arms in tandem. These are known as \textbf{ganged potentiometers}.
            \bigbreak \noindent 
            Finally, unlike fixed resistors, pots are available in limited sizes. Typical values use multiples of 1 or 5 for the total resistance
            \bigbreak \noindent 
            The output voltage at the wiper of a potentiometer changes because the wiper splits the total resistive track into two smaller resistor sections, and the voltage drops across each section in direct proportion to its length
        \item \textbf{Rheostat}: If only the wiper and only one of the two end terminals is used, then the device is referred to as a rheostat. Rheostats are often used to limit current. The value of the rheostat ranges from 0 up to the nominal value of the potentiometer.
            \bigbreak \noindent 
            \fig{.6}{./figures/16.png}
        




    \end{itemize}

    \pagebreak 
    \unsect{Parallel resistive circuits}
    \begin{itemize}
        \item \textbf{The parallel connection}:  In a parallel connection the voltage is the same across each component. 
            \bigbreak \noindent 
            \fig{.6}{./figures/17.png}
            \bigbreak \noindent 
            Although it may be drawn oddly, there are only two common connection nodes in this configuration: one along the top and the other along the bottom
            \bigbreak \noindent 
            If we were to take a voltmeter and place it across any of these blocks, the meter would read the same voltage across each because the probes are always contacting the same two nodes
            \bigbreak \noindent 
            \fig{.6}{./figures/18.png}
            \bigbreak \noindent 
            In this layout some elements are in series and some are in parallel. Only elements D and E are strictly in parallel here as they are the only two that connect to the same two nodes, and thus must see the same voltage.
            \bigbreak \noindent 
            Items C and F might at first appear to be in parallel with them but upon closer inspection it should be noted that C and F are in series with each other. That is, C and F will split the voltage that D or E sees, in accordance with the voltage divider rule. In other words, it is not true that the voltage across C or across F must be the same as the voltage across D or E
        \item \textbf{Combining Parallel Components}: Our first step is to determine how to combine parallel components in order to create a single equivalent component. Unlike series connections, this can be a little more time consuming.
        \item \textbf{Sources in parallel}: First, voltage sources are not placed in parallel as a general rule
            \bigbreak \noindent 
            \fig{.6}{./figures/19.png}
            \bigbreak \noindent 
            The reason is because a parallel connection requires the same voltage across each component. This would be impossible to achieve with each source trying to maintain a different voltage across the same two nodes. This may result in damage to the sources (e.g., exploding batteries). The main exception to this rule is if the sources have the same potential and the goal is to extend operational life (e.g., multiple D cell batteries in parallel).
        \item \textbf{Current sources in parallel}: When it comes to current sources, they may simply be added together, however, like series voltage sources, polarity is important. Referring to the three parallel current sources in Figure 4.4, the middle and right sources are pumping current into the top node while the left source is feeding the bottom node. You can also think of the left source as draining the top node. Thus, these three sources together behave like a single source of five amps feeding the top node (eight amps in, three amps out). 
            \bigbreak \noindent 
            \fig{.6}{./figures/20.png}
        \item \textbf{Resistors in Parallel}: When placed in parallel, resistor values do not add they way they do in series connections.
            \bigbreak \noindent 
            \fig{.6}{./figures/21.png}
            \bigbreak \noindent 
            Recall that 
            \begin{align*}
                R = \frac{\rho l}{A}
            .\end{align*}
            Notice that two of the same resistor in parallel would be double the area, which resistivity and length unchanged. Thus, the resistance is halved. So, we see that placing resistors in parallel results in a decrease in net resistance, the opposite of the series case.
            \bigbreak \noindent 
            Recall that \textbf{conductance} $G$ is given by the reciprocal of resistance
            \begin{align*}
                G = \frac{A}{\rho l}
            .\end{align*}
            Now we can see that the increase in area creates a proportional increase in conductivity. If we generalize this for $N$ resistors we find that the equivalent conductance of a group of parallel resistors is their sum
            \begin{align*}
                G_{\text{total}} = G_{1} + \dots + G_{N}
            .\end{align*}
            Thus,
            \begin{align*}
                \frac{1}{R_{\text{parallel}}} = \frac{1}{R_{1}} + \dots + \frac{1}{R_{N}}
            .\end{align*}
            Hence,
            \begin{align*}
                R_{\text{parallel}} = \frac{1}{\frac{1}{R_{1}} + \dots + \frac{1}{R_{N}}}
            .\end{align*}
            For $N$ resistors of equal resistance, 
            \begin{align*}
                R_{\text{parallel}} = \frac{1}{\frac{N}{R}} = \frac{R}{N}
            .\end{align*}
            Also, consider two parallel resistors $R_{1}$ and $R_{2}$, with $R_{1} < R_{2}$. Then,
            \begin{align*}
                R_{P} = \frac{1}{\frac{1}{R_{1}} + \frac{1}{R_{2}}}
            .\end{align*}
            Consider $R_{1}$ alone,
            \begin{align*}
                R_{P} = \frac{1}{\frac{1}{R_{1}}} = R_{1}
            .\end{align*}
            Now, consider what happens when we add $R_{2}$. Since resistance is non-negative, adding $1/R_{2}$ admits 
            \begin{align*}
                \frac{1}{R_{1}} + \frac{1}{R_{2}} > \frac{1}{R_{1}}
            .\end{align*}
            Thus,
            \begin{align*}
                \frac{1}{\frac{1}{R_{1}} + \frac{1}{R_{2}}} < \frac{1}{\frac{1}{R_{1}}} = R_{1}
            .\end{align*}
            Thus, the combined resistance is smaller than the smallest resistor in the parallel series.
        \item \textbf{Product-Sum Rule}: For two resistors, 
            \begin{align*}
                R_{P} = \frac{R_{1}R_{2}}{R_{1} + R_{2}}
            .\end{align*}
            It is only applicable to two resistors, however, it can be used repeatedly on pairs of resistors within larger groups instead of using
            \bigbreak \noindent 
            Now, suppose we write
            \begin{align*}
                N = \frac{R_{2}}{R_{1}}
            ,\end{align*}
            so that
            \begin{align*}
                R_{2} = \frac{N}{R_{1}}
            .\end{align*}
            Now,
            \begin{align*}
               R_{P} = \frac{R_{1}NR_{1}}{R_{1} + NR_{1}}  = \frac{N}{N+1}R_{1}
            .\end{align*}
            This is useful for quick estimates of parallel resistor pairs. If we have a resistor ratio of 3:1, then the equivalent would be $3/4$ths the value of the smaller resistor. For a $24 k\Omega$ and a $8 k \Omega$, the equivalent would be a $6 k \Omega$ resistor.
        \item \textbf{Kirchhoff's Current Law (KCL)}: Just as Kirchhoff's voltage law is a key element in understanding series circuits, Kirchhoff's current law (KCL) is the operative rule for parallel circuits. It states that the sum of all currents entering and exiting a node must sum to zero. Alternately, it can be stated as the sum of currents entering a node must equal the sum of currents exiting that node. As a pseudo formula: 
            \begin{align*}
                \sum I \leftarrow = \sum I \rightarrow
            .\end{align*}
            Recall that a node is a connection area wherein the voltage is the same
            \bigbreak \noindent 
            \fig{.6}{./figures/22.png}
            \bigbreak \noindent 
            We see that there are two nodes in a parallel configuration. For each of these nodes, whatever currents flow in must be balanced by the same total current flowing out.
            \bigbreak \noindent 
            There is a current source feeding the top node. At the three-way connection with $R_{1}$ and $R_{2}$, KCL along with Ohm's law dictate that the incoming current from the source must split with one portion exiting by flowing down through $R_{1}$ and the remainder exiting by flowing down through $R_{2}$.
            \bigbreak \noindent 
            At the bottom node (ground) these two resistor currents must combine and equal the current that is being drawn back into the current source, which must be the same as the current originally produced by the source. This is perfectly sensible because it would be impossible for the current source to establish two different currents; one at its input and the other at its output. To assert otherwise would be a little like saying you're two different heights depending on whether you measure from your toes up to your head versus from your head down to your toes. Thus, a simplistic way of stating KCL is “What goes in must come out”
        \item \textbf{The Current Divider Rule (CDR)}: Just as series circuits follow the voltage divider rule (voltage dividing in proportion to resistance), parallel circuits follow the current divider rule which states that current divides in reverse proportion to resistance (i.e., in direct proportion to conductance).
            \bigbreak \noindent 
            First, we note that the voltage across the pair must equal the entering current times the effective resistance of the pair.
            \begin{align*}
                V = I_{T}\frac{R_{1}R_{2}}{R_{1} + R_{2}}
            .\end{align*}
            The current through each of the resistors must be this voltage divided by the corresponding resistance. For the current through $R_{1}$
            \begin{align*}
                I_{R_{1}} = \frac{V}{R_{1}} = I_{T}\frac{R_{2}}{R_{1} + R_{2}}
            .\end{align*}
            Similarly, for the current through $R_{2}$:
            \begin{align*}
                I_{R_{2}} = I_{T}\frac{R_{1}}{R_{1} + R_{2}}
            .\end{align*}
            Hence,
            \begin{align*}
                I_{X} = I_{T} \frac{R_{Y}}{R_{X} + R_{Y}}
            .\end{align*}
            This rule is convenient in that you don't have to compute the parallel equivalent resistance, but remember, it is valid only when there are just two resistors involved.
        \item \textbf{Fuses and circuit breakers}: Most electrical and electronic devices are designed to run off of a constant voltage source As such, parallel connections are ideal when a variable number of devices are to be powered because, unlike a series connection, each device will see the same voltage rather than having that voltage split among them.
            \bigbreak \noindent 
            There is another advantage to parallel connection. Consider
            \bigbreak \noindent 
            \fig{.6}{./figures/23.png}
            \bigbreak \noindent 
            Typically, when a lamp fails, it fails creating an open. If any of the lamps in the series connection fails, current in the loop is interrupted and all of the lamps go out. In contrast, if one of the lamps in the parallel connection fails, the remaining lights will stay on. Further, in the series connection, there is no immediate way to determine which lamp has failed as none of them will be lit. Each lamp has to be tested in turn to determine the faulty unit. In the parallel system, it is obvious which lamp has failed because that's the one that's out.
            \bigbreak \noindent 
            In a home, school or commercial setting, any number of devices may be plugged into wall sockets with the expectation that each unit will receive the same voltage and operate correctly due to the parallel nature of the system. As convenient as this is, there is a practical problem. Every device that is added to the system presents another path for current flow. Due to KCL, this means that the total current drawn from the source must increase. This creates a potential problem because if too much current is drawn, there can be an appreciable voltage drop across the wires feeding the outlets as they can no longer be assumed to be zero ohms. This lowered voltage may present a problem for the attached devices, but more importantly, the large current results in unexpected heating of the wires because of the power dissipation in them ($I^{2}R$) and creates a possible fire hazard. Consequently, we must include some method to limit the current to a safe maximum value.
            \bigbreak \noindent 
            Adding another load in a parallel circuit increases the total current drawn from the source because it creates a new, separate path for electricity to flow.  Each load you add creates an independent branch connected across the same voltage source.
            \bigbreak \noindent 
            The power source provides the same push (voltage) to every branch. Each new load draws its own separate share of current based on its own resistance
            \bigbreak \noindent 
            According to Kirchhoff's Current Law, the total current leaving the source is the sum of the currents going through each individual branch. When you add another branch current to the total, the overall sum goes up
            \bigbreak \noindent 
            Because there are more pathways for the charge to move through, the overall or equivalent resistance of the entire circuit actually decreases.
            \bigbreak \noindent 
            Since the resistance decreases, 
            \begin{align*}
                I = \frac{V}{R}
            \end{align*}
            increases. The common approach is a fuse or circuit breaker. These devices are placed between the voltage source and the various loads, and see their combined current. If this current becomes too great, the device interrupts the current flow by opening the circuit. Nothing will operate but nothing catches fire, either.
            \bigbreak \noindent 
            Fuses are used typically in electronic systems, motors, automotive subsystems and the like. Whether they're the cylindrical type found in electronic instrumentation or the blade type common in automobiles, fuses are one-shot devices. They consist of a fine wire or metal link that will heat up and melt if it experiences too high of a current. Upon melting, circuit continuity is lost and current flow stops. This is called a “blown” fuse. The fuse will then need to be replaced once the problem has been fixed (i.e., whatever it was that caused excessive current to flow in the first place).
            \bigbreak \noindent 
            Fuses have a nominal current rating, for example, two amps. This does not mean that as soon as two amps is reached, the fuse blows. It will take some time for the fuse to heat up. The higher the current is above the nominal rating, the faster the fuse will blow.
            \bigbreak \noindent 
            \fig{.6}{./figures/25.png}
            \bigbreak \noindent 
            \fig{.6}{./figures/24.png}
            \bigbreak \noindent 
            Fuses are also available in fast-blow and slow-blow forms. Fast-blow fuses will activate many times faster than standard fuses while slow-blow types will require much long to activate (this prevents the fuse from activating accidentally with certain kinds of loads that exhibit a high initial startup current but that falls back to a more modest level once running smoothly, such as a motor).
            \bigbreak \noindent 
            Fuses are generally inexpensive and effective. The obvious downside to a fuse is that they need to be replaced once they do their job. A circuit breaker solves this problem. 
            \bigbreak \noindent 
            A circuit breaker, or just breaker for short, can be thought of as a resettable fuse. Breakers do not have a melting link that interrupts current flow. Instead, a breaker operates like an intelligent switch: if the current is too large, the switch is thrown open. The operator can then fix the problem and reset the breaker (i.e., return the switch to the operating position).
            \bigbreak \noindent 
            \fig{.6}{./figures/26.png}





    \end{itemize}

    \unsect{Series-Parallel Resistive Circuits}
    \begin{itemize}
        \item \textbf{Intro}: The key to analyzing basic series-parallel circuits is in recognizing portions of the circuit, or sub-circuits, that exhibit a series or parallel configuration by themselves, and then applying the series and parallel analysis rules to those sections. Ohm's law, KVL and KCL may be used in turn to “chip away” at the problem until all currents and voltages are found. As individual voltages and currents are determined, this makes it easier to apply these rules to determine other values. Given this observation, the number of potential solution paths tends to grow exponentially as the number of components increases. As a consequence, when faced with the same circuit, six people may solve it six different ways, no particular way being more or less correct than any other. The only thing we can say is that some solution paths might be more computationally efficient than others, meaning they take less work 
        \item \textbf{Series-Parallel Connections}: 
            \bigbreak \noindent 
            \fig{.6}{./figures/27.png}
            \bigbreak \noindent 
            It should be obvious that this system is neither a simple series connection nor a simple parallel connection, but rather some combination of the two. In order for it to be series, the current through each component would have to be the same. This does not have to be true here. Consider the connection node directly above block C. This also connects to items B, D and E. All we really know about that node is that there are four paths leading to it and that the four associated currents must sum to zero, thanks to KCL. There is nothing that requires the current through, say, block D to be the same as that through block E. Similarly, there is nothing that stipulates that the voltage across block C must be the same as that across block F. After all, that is the hallmark of a parallel configuration, and that would require those two blocks to be connected to the same two nodes, which they aren't.
            \bigbreak \noindent 
            Notice some of the components are in series among themselves, and some are in parallel among themselves. We might recognize these groupings as sub-circuits.
            \bigbreak \noindent 
            For example, it is true that blocks A and B are in series with each other because the current through one of them must be the same as the current through the other. The same is true of blocks C and F; they are in series with each other. Similarly, we can say that blocks D and E are in parallel with each other because they must see the same voltage as they are each connected to the same two nodes
        \item \textbf{Ladders}: A ladder is a unique series-parallel configuration. It is arranged in a cascade of series and parallel connections.
            \bigbreak \noindent 
            \fig{.6}{./figures/28.png}
            \bigbreak \noindent 
            The main thing to note is that the ladder network consists of “sections loading sections” repeatedly. By loading, we mean that this element (the load) draws off current and power from the prior section. As a result, simple two-resistor voltage dividers cannot be used.
            \bigbreak \noindent 
            The only proper voltage divider in this configuration is between $R_{7}$ and $R_{8}$. The reason for this becomes apparent if we imagine determining the equivalent resistance of the network by placing an ohmmeter at terminals $a$ and $f$. The ohmmeter will apply a current to node $a$ which passes through $R_1$. At node $b$ this current splits, part going down $R_2$ and part continuing through $R_3$. Clearly then, $R_1$ and $R_2$ are not in series as they do not see the same current. Further, the current that flows through $R_3$ enters node c where it splits again, part down $R_4$ with the remainder flowing through $R_5$. Thus, $R_3$ and $R_4$ are not in series, either. The same sort of thing happens again at node $d$. This would continue for as many rungs as are on the ladder, the sole exception being the final two resistors which are in series ($R_7$ and $R_8$ here).
            \bigbreak \noindent 
            How then do we find the equivalent resistance of this network? We begin at the end farthest from the open terminals. 
            We have already noted that \(R_7\) and \(R_8\) are in series. This pair,
            if treated as a single resistance, is in parallel with \(R_6\), or
            \[
                R_6 \parallel (R_7 + R_8).
            \]
            This group of three is in series with \(R_5\), yielding
            \[
                R_5 + \bigl(R_6 \parallel (R_7 + R_8)\bigr).
            \]
            This new group of four is in parallel with \(R_4\), yielding
            \[
                R_4 \parallel \bigl(R_5 + (R_6 \parallel (R_7 + R_8))\bigr).
            \]
            This group is in series with \(R_3\), and that new group is in turn in
            parallel with \(R_2\). Finally, that penultimate grouping is in series
            with \(R_1\). Therefore, the equivalent resistance of the network is
            \[
                R_{\mathrm{eq}} = R_1 + \left[ R_2 \parallel \left( R_3 + \left[ R_4 \parallel \left( R_5 + \left[ R_6 \parallel (R_7 + R_8) \right] \right) \right] \right) \right].
            \]
            Surely, it is a tad tedious to compute the value, but it is not
            particularly difficult. A word to the wise: if the goal is to determine
            the voltages at the various nodes or the myriad branch currents, it will
            be helpful to keep track of the equivalent resistances of the various
            groupings.
            \bigbreak \noindent 
            For example, if \(V_{cf}\) is known and the goal is to find \(V_{df}\),
            the voltage-divider relationship would be
            \[
                V_{df}
                =
                V_{cf}
                \frac{R_6 \parallel (R_7 + R_8)}
                {R_5 + \bigl(R_6 \parallel (R_7 + R_8)\bigr)}.
            \]
            Note that || has higher precedence than + or -, therefore || is performed first, rather like multiplication when mixed with addition and subtraction.
        \item \textbf{Any Route Works}: To find some voltage $V_{XY}$, start at point $X$ and proceed along any convenient route to $Y$, subtracting voltage rises (polarities of - to +) and adding voltage drops (polarities of + to -) along the way.
        \item \textbf{R-2R Ladder}: As promised earlier in the chapter, it's time to look at a ladder network. While ladders can use any resistance values desired, certain ratios turn out to be eminently practical. Of particular note is the R-2R ladder
            \bigbreak \noindent 
            In this configuration, only two different resistor values are used; some initial value and a second value of precisely twice that size
            \bigbreak \noindent 
            \fig{.6}{./figures/29.png}
            \bigbreak \noindent 
            It turns out that this arrangement offers something very unique: a halving of successive node voltages and branch currents. This is extremely useful in digital to analog conversion circuits; the sort of hardware that would turn the digital bits of an MP3 or WAV file into listenable analog signals to feed loudspeakers or headphones.
            \bigbreak \noindent 
            Starting at the far right end, it should be obvious that the last two 6 $k \Omega$ resistors comprise a 50\% voltage divider, and thus $V_{e}$ will half of $V_{d}$ 
            \bigbreak \noindent 
            The combination of the final three resistors works out to 12 $k \Omega$ || (6 $k \Omega$ + 6 $k \Omega $), or 6 $k \Omega$ again. Thus, the divider from c to d is also 50\%. This process repeats as we move to the left, each equivalent working out to $R$ and resulting in a 50\% voltage divider. It will not matter what value is chosen for $R$. As long as $2R$ is used for the other set of resistors, the result will always be successive 50\% voltage dividers.
        \item \textbf{Bridges}: A bridge is a particular series-parallel configuration consisting of two pairs of series connected elements placed in parallel. 
            \bigbreak \noindent 
            \fig{.6}{./figures/30.png}
            \bigbreak \noindent 
            A typical use for a bridge is the measurement of some environmental quantity. There are certain resistive devices that are sensitive to environmental change. Examples include photoresistors, whose resistance is a function of light level; and thermistors, whose resistance is a function of temperature.
            \bigbreak \noindent 
            To understand the operation, the voltages $V_{b}$ and $V_{c}$ are established by the voltage dividers comprised of $R_1$ and $R_2$; and $R_3$ and $R_4$, respectively. With ordinary resistors, these voltages would be unchanging and therefore $V_{bc}$ would also be a fixed value. We shall set it up such that $V_{b}$ and $V_{c}$ are the same voltage, and thus $V_{bc}$ will be zero. Now assume that we replace $R_1$ with a photoresistor. As the light level increases, its resistance decreases. This will cause $V_{b}$ to rise which will cause $V_{bc}$ to go positive. If the light level decreases, the resistance of $R_1$ will increase, forcing Vb to drop which will cause $V_{bc}$ to go negative. The greater the change in light, the larger $V_{bc}$ will be, and its sign indicates whether the light level has increased or decreased compared to the original set point. By simply placing some manner of voltmeter between points $b$ and $c$, we can create a display that indicates light levels.
            \bigbreak \noindent 
            If we want to flip the sign, we would place the photoresistor in the position of $R_2$ instead of $R_1$. Further, it is possible to increase the sensitivity by using sensors at opposite corners (e.g., $R_1$ and $R_4$), and comparative or differential measurements are possible by using both sides (e.g., $R_1$ with $R_3$).


    \end{itemize}

    \pagebreak 
    \unsect{Analysis theorems and techniques}
    \begin{itemize}
        \item \textbf{Intro}: We begin by considering a more realistic model for constant voltage and constant current sources. The ideal voltage source produces a stated potential forever, and without regard to what it is connected to. The ideal current source behaves similarly; it will always produce the same current regardless of its load. These expectations are not realistic. For example, if we were to place a solid bar of copper across the terminals of a voltage source, the bar may exhibit a resistance of mere milliohms, implying an output current of thousands of amps. Similarly, if we were to disconnect a current source from any load, its effective load would be the resistance of the air between its terminals and Ohm's law would dictate an output voltage of perhaps thousands or even millions of volts. Real world sources do not behave like this.          
            \bigbreak \noindent 
            Thus, ideal voltage and current sources cannot exist in the real world because they would theoretically produce infinite or destructive amounts of power.
        \item \textbf{Internal resistance (Voltage source)}: A highly precise model for any voltage source or current source could be fairly complex but for general purpose work, we can greatly improve our ideal sources by simply adding a resistor to them. This resistance is referred to as the internal resistance of the source.
            \bigbreak \noindent 
            It is important to understand that this is not a resistor, as in an internal component that can be altered, but rather a mathematical addition to the source that better predicts how it will behave. Further, the value of this effective resistance can be deduced in a laboratory via proper measurements.
            \bigbreak \noindent 
            The model for a voltage source adds a resistance in series
            \bigbreak \noindent 
            \fig{.6}{./figures/31.png}
            \bigbreak \noindent 
            This resistance sets an upper limit on the source's current output. Even if the output terminals are shorted, the maximum current will be dictated via Ohm's law to be the source voltage divided by the internal resistance, or E/R. Obviously, this internal resistance will create some voltage divider effect with the attached load. To minimize this effect, the internal resistance should be as small as practicably possible. Thus,
            \begin{quote}
               "\textit{The ideal internal resistance of a voltage source is zero ohms (a short). }" 
            \end{quote}
            \bigbreak \noindent 
            Given a zero ohm internal resistance, this improved model reverts back to the original ideal source. In the case of a lab power supply, the internal resistance typically will be a small fraction of an ohm.
            \bigbreak \noindent 
            An ideal voltage source has zero internal resistance, meaning it maintains a completely constant terminal voltage regardless of how much current is drawn by the connected circuit. It provides a fixed voltage output at all times.
            \bigbreak \noindent 
            The voltage drop is then
            \begin{align*}
                V = IR = 0 
            .\end{align*}
        \item \textbf{Internal resistance (Current source)}: For a current source, the improved model adds a resistance in parallel
            \bigbreak \noindent 
            \fig{.6}{./figures/32.png}
            \bigbreak \noindent 
            This resistance sets an upper limit on the source's voltage output. If the output terminals are opened, the maximum voltage will no longer produce a huge voltage. Instead, it is dictated by Ohm's law to be the source current times the internal resistance, or $IR$. This internal resistance will create some current divider effect with the attached load. To minimize this effect, the internal resistance should be as large as practicably possible. Thus,
            \begin{quote}
               "\textit{The ideal internal resistance of a current source is infinite ohms. }" 
            \end{quote}
            \bigbreak \noindent 
            Given an infinite ohm internal resistance (i.e., an open), this improved model reverts back to the original ideal source. From here on, whenever we deal with practical voltage and current sources, we understand that these sources have some associated internal resistance, even if they're not shown explicitly in a schematic diagram. Further, whenever we talk about ideal sources, we simply use a short for the internal resistance of a voltage source and an open for the internal resistance of a current source.
            \bigbreak \noindent 
            An ideal current source has infinite internal resistance, meaning it delivers a completely constant current regardless of the voltage across its terminals.
            \bigbreak \noindent 
            When you connect an external load (\[R_{load}\]): 
            \begin{itemize}
                \item The total current splits into two paths: some goes through \[R_{internal}\] (bleed current), and the rest goes to your load.
                \item Because they are in parallel, both paths experience the exact same voltage (V)
            \end{itemize}
            Since
            \begin{align*}
                I_{\text{internal}} = \frac{V}{\infty} = 0
            ,\end{align*}
            there is no internal current bleed.
        \item \textbf{Open load}: An open load in a circuit is a condition where the intended device or component (the "load") is disconnected, broken, or missing, creating a break in the electrical pathway that stops current from flowing
            \bigbreak \noindent 
            Because the conductive path is incomplete, the electrical current drops to zero amperes ($I \approx  0$).
            \bigbreak \noindent 
            The gap or break acts as an extremely high (effectively infinite) resistance barrier.
            \bigbreak \noindent 
            While current cannot flow, the full source voltage can still be measured across the break or the terminals where the load should be
            \bigbreak \noindent 
            To continue, if we look at the open load case, for the voltage source the load current would be zero and the load voltage would be the entire source voltage of $E$. For the current source, the load would also see no current and its voltage would be the voltage appearing across its internal resistance which is R times the current $E/R$, or just $E$. Thus, the two behave identically at the load limits.
        \item \textbf{Source Equivalences}: For any simple voltage source consisting of an ideal voltage source with a series internal resistance, an equivalent current source may be created. Similarly, for any simple current source consisting of an ideal current source with a parallel internal resistance, an equivalent voltage source may be created. By “equivalent”, we mean the load currents must be identical for both circuits given any value of load resistance. (Note that if the currents are the same then the voltages must also be the same due to Ohm’s Law.)
            \bigbreak \noindent 
            Given a voltage source, the maximum current is developed when the load is shorted producing a current of $E/R$. Under that same load condition, all of the current from the current source version must be flowing through the load. Therefore, the value of the equivalent current source must be the maximum current of $E/R$. It would be nonsensical to use a current source that was larger or smaller than this value.
            \bigbreak \noindent 
            To continue, if we look at the open load case, for the voltage source the load current would be zero and the load voltage would be the entire source voltage of $E$. For the current source, the load would also see no current and its voltage would be the voltage appearing across its internal resistance which is $R$ times the current $E/R$, or just $E$. Thus, the two behave identically at the load limits.
            \bigbreak \noindent 
            Similarly, if we start with a current source, an open load produces the maximum load voltage of $IR$. Therefore, the equivalent voltage source must have a value of $IR$. For the current source, a shorted load would produce a load current equal to the source value, or $I$. The voltage source version would produce a current of $E/R$, where the value of $E$ was just found to be equal to $IR$, and thus the load current would be $IR/R$, or just $I$. Once again, the two versions behave identically at the load limits
            \bigbreak \noindent 
            To summarize the process of source conversion:
            \begin{itemize}
                \item The internal resistance will be the same for both versions.
                \item If converting from a voltage source to a current source, the value of the current source will be the maximum current available from the voltage source (shorted load case), and is equal to $E/R$.
                \item If converting from a current source to a voltage source, the value of the voltage source will be the maximum voltage available from the current source (opened load case), and is equal to $IR$.
            \end{itemize}
            If a multi-source is being converted (i.e., voltage sources in series or current sources in parallel), first combine the sources to arrive at the simplest source and then do the conversion. Do not convert the sources first and then combine them as you will wind up with series-parallel configurations rather than simple sources.
            \bigbreak \noindent 
            Judicious use of source conversions can sometimes simplify multi-source circuits by allowing converted sources to be combined, resulting in a single source
            \bigbreak \noindent 
            Now, consider
            \bigbreak \noindent 
            \fig{.6}{./figures/33.png}
            \bigbreak \noindent 
            This circuit is unlike any circuit we have seen so far. Although we have analyzed circuits using multiple voltage sources, they have always been in a simple series loop. As such, their voltages may be added together to find a single equivalent voltage source. Such is not the case here. In this circuit the voltage sources are part of a series-parallel network and thus their potentials cannot be merely added. In fact, there are no further simplifications to be performed on this circuit using basic series-parallel techniques. We appear to be stuck.
            \bigbreak \noindent 
            But we're not. This circuit can be simplified into a straight all-parallel network through the use of source conversions. The $E_{1}$, $R_{1}$ combo can be converted into one current source while the $E_{2}$, $R_{2}$ combo can be converted into a second source. Once these are converted, the resulting circuit will consist of two current sources and three resistors, all in parallel
            \bigbreak \noindent 
            Consider
            \bigbreak \noindent 
            \fig{.6}{./figures/34.png}
            \bigbreak \noindent 
            The converted circuit is then
            \bigbreak \noindent 
            \fig{.6}{./figures/35.png}
            \bigbreak \noindent 
            The equivalent resistance is 
            \begin{align*}
                V_{\text{EQ}} = 1000 || 4000 || 5000 = 689.7 \Omega
            .\end{align*}
            Since the total current is $15 + 1.5 = 16.5$mA,
            \begin{align*}
                V_{b} = 0.0165(689.7) = 11.38V
            .\end{align*}
            We can now take this voltage and apply it back to the original (unconverted) circuit. With this potential known, it is relatively easy to determine the other currents and voltages using KVL and Ohm's law.
            \bigbreak \noindent 
            Looking at the \(1\,\mathrm{k}\Omega\) resistor, the voltage across it is
            \[
                15\,\mathrm{V} - 11.38\,\mathrm{V} = 3.62\,\mathrm{V}.
            \]
            Therefore, the current through it is
            \[
                I_1 = \frac{3.62\,\mathrm{V}}{1\,\mathrm{k}\Omega}
                = 3.62\,\mathrm{mA}.
            \]
            Similarly, the voltage across the \(4\,\mathrm{k}\Omega\) resistor is
            \[
                11.38\,\mathrm{V} - 6\,\mathrm{V} = 5.38\,\mathrm{V},
            \]
            which yields a current of
            \[
                I_2 = \frac{5.38\,\mathrm{V}}{4\,\mathrm{k}\Omega}
                = 1.345\,\mathrm{mA}.
            \]
            Both currents flow from left to right. The third current, which flows downward through the \(5\,\mathrm{k}\Omega\) resistor, is
            \[
                I_3 = \frac{11.38\,\mathrm{V}}{5\,\mathrm{k}\Omega}
                = 2.276\,\mathrm{mA}.
            \]
            According to KCL, the current entering node \(b\) must equal the total current leaving it. The entering current is \(3.62\,\mathrm{mA}\), while the sum of the exiting currents is
            \[
                I_2 + I_3
                = 1.345\,\mathrm{mA} + 2.276\,\mathrm{mA}
                = 3.621\,\mathrm{mA}
                \approx 3.62\,\mathrm{mA}.
            \]
            Therefore, KCL is satisfied within the rounding precision of the calculated values.
            \bigbreak \noindent 
            \textbf{Note:} In the equivalent circuit, nodes $a$ and $c$ disappeared. This brings up an important point. Equivalent circuits are equivalent in terms that the items connected to the equivalent behave the same way as they do with the original circuit. That does not mean that the items within the equivalent see the same current or voltage. We do not expect the voltage across the 1 $k \Omega$ or 4 $k \Omega$ resistors in the converted version to be the same as those in the original version.
        \item \textbf{Linear and bilateral circuits}: A linear circuit has an output that is directly proportional to its input, while a bilateral circuit allows current to flow and behave the same way in either direction
        \item \textbf{Linear circuits}: A circuit where values like resistance, inductance, and capacitance stay constant and do not change when voltage or current changes.  They follow Ohm's law and obey the principle of superposition 
        \item \textbf{Bilateral circuits}: A network or circuit whose properties and current-voltage behavior are identical regardless of the direction of the current. Reversing the voltage polarity reverses the current by the exact same magnitude, meaning the circuit is bidirectional.
        \item \textbf{Recall multiple voltage sources}: A voltage source acts like a pump for electrical charge, with a positive terminal pushing current out one way and a negative terminal pulling it back
            \bigbreak \noindent 
            When you connect the positive terminal of one battery to the negative terminal of another (+ to -), both sources pump charge in the same direction. Their efforts combine, so you add their voltages
            \bigbreak \noindent 
            When you connect opposing terminals together (+ to + or - to -), the two sources try to force current in opposite directions. One source tries to raise the potential while the other tries to lower it.
        \item \textbf{Short}: A short means a direct connection with zero ohms of resistance between two points.
        \item \textbf{Open}: An open circuit is an incomplete electrical path where the connection is broken, resulting in infinite resistance and zero current flow.
        \item \textbf{Superposition Theorem}: The superposition principle states that the total voltage or current at any point in a linear circuit with multiple independent sources is the algebraic sum of the individual voltages or currents produced by each source acting alone
            \bigbreak \noindent 
            Superposition allows the analysis of multi-source series-parallel circuits. Superposition can only be applied to networks that are linear and bilateral. Further, it cannot be used to find values for non-linear functions, such as power, directly. Fortunately, if the circuit contains nothing but resistors, and ordinary voltage sources and current sources, the circuit will be a linear bilateral network. Also, although power is a square law function (i.e., it is proportional to the square of voltage or current), it can be computed from the resulting voltage or current values so this presents no limits to analysis.
            \bigbreak \noindent 
            The basic idea is to determine the contribution of each source by itself, and then adding the results to get the final answer(s). Consider
            \fig{.6}{./figures/36.png}
            \bigbreak \noindent 
            Here we have two voltage sources driving a pair of resistors. It is a basic series circuit. As the two sources are in opposition, the net voltage is 10.5 volts while the total resistance is 500 $\Omega$. This yields a circulating current of 21 mA. 
            \bigbreak \noindent 
            From earlier work we have discovered that the ideal internal resistance of a voltage source is a short, or zero ohms. Thus, we wind up with two circuits: one with a 12 volt source that drives 100 $\Omega$, 400 $\Omega$ and a short; and a second circuit with a 1.5 volt source that drives pretty much the same thing, but with an opposing current direction.
            \bigbreak \noindent 
            The first circuit produces a clockwise current of $12 V/ 500 \Omega$, or 24 milliamps. Meanwhile, the second circuit produces a counterclockwise current of $1.5 V / 500 \Omega$, or 3 mA. As these two currents oppose each other, the resulting current is $24$ mA $- 3$ mA, or 21 mA, the same value originally computed
            \begin{itemize}
                \item For every voltage or current source in the original circuit, create a new subcircuit. The sub-circuits will be identical to the original except that all sources other than the one under consideration will be replaced by their ideal internal resistance. This means that all remaining voltage sources will be shorted and all remaining current sources will be opened.
                \item Label the current directions and voltage polarities on each of the new subcircuits, as generated by the source under consideration.
                \item Solve each of the sub-circuits for the desired voltages and/or currents using standard series-parallel analysis techniques. Make sure to note the voltage polarities and current directions for these items.
                \item Add all of the contributions from each of the sub-circuits to arrive at the final values, being sure to account for current directions and voltage polarities in the process.
            \end{itemize}
            Reconsider
            \bigbreak \noindent 
            \fig{.6}{./figures/34.png}
            \bigbreak \noindent 
            The two sub circuits are
            \bigbreak \noindent 
            \fig{.6}{./figures/37.png}
            \bigbreak \noindent 
            \fig{.6}{./figures/38.png}
            \bigbreak \noindent 
            Solving for $V_{b}$ on the first yields $V_{b} = 10.34V$, solving for $V_{b}$ on the second yields $V_{b} = 1.034V$. By the superposition principle, we add both contributions. Thus, $V_{b}$ in the original circuit is approximately
            \begin{align*}
                V_{b} \approx 11.37V
            .\end{align*}
            If the $6V$ source was reversed, it would pull node $b$ toward $-6V$. Its contribution would then be negative, since $V_{b}$ in this case would be below ground.
        \item \textbf{Ports}: In electrical circuit theory, a port is a pair of terminals connecting a circuit or network to the outside world where the current entering one terminal must equal the current exiting the other terminal
            \bigbreak \noindent 
            The net current flowing into the pair of terminals must be zero. This means the current entering one terminal is equal and opposite to the current leaving the companion terminal
        \item \textbf{Thevenin's Theorem}: Any single port linear network can be reduced to a simple voltage source, $E_{th}$, in series with an internal resistance, $R_{th}$.
            \bigbreak \noindent 
            The phrase “single port network” means that the circuit is cut in such a way that only two connections exist to the remainder of the circuit (two connection points makes up one port). The remainder of the circuit may be a single component or a large multi-component sub-circuit. As there are many ways to cut a typical circuit, there are many possible Thevenin equivalents. The important thing is that there are only two connection points between the two portions of the circuit and neither point has to be ground.
            \bigbreak \noindent 
            Consider
            \bigbreak \noindent 
            \fig{.6}{./figures/39.png}
            \bigbreak \noindent 
            Suppose we cut the circuit immediately to the left of $R_4$. That is, we will find the Thevenin equivalent that drives $R_4$. The first step is to make the cut, removing the remainder of the circuit (in this case, just $R_4$). We then determine the open circuit output voltage. This is the maximum voltage that could appear between the cut points and is called the Thevenin voltage, $E_{th}$.
            \bigbreak \noindent 
            In a circuit such as this, basic series-parallel analysis may be used to find $E_{th}$. This process turns out to be quite straightforward in this particular circuit. Due to the open, no current flows through $R_3$, thus no voltage is developed across $R_3$, and therefore $E_{th}$ must equal the voltage developed across $R_2$ which may be obtained via a voltage divider with resistor $R_1$ and source $E$.
            \bigbreak \noindent 
            \fig{.6}{./figures/40.png}
            \bigbreak \noindent 
            The second part is finding the Thevenin resistance, $R_{th}$. Beginning with the “cut” circuit, replace all sources with their ideal internal resistance (thus shorting voltage sources and opening current sources). From the perspective of the cut point, look back into the circuit and simplify by making appropriate series and parallel combinations to determine its equivalent resistance.
            \bigbreak \noindent 
            \fig{.6}{./figures/41.png}
            \bigbreak \noindent 
            We find that $R_1$ and $R_2$ are in parallel, and this combination is then in series with $R_3$. Thus, $R_{th} = R_3 + (R_1 || R_2)$. A common error is to determine the equivalent resistance that the source drives. Remember, everything is determined from the vantage point of the cut or port.
            \bigbreak \noindent 
            Thevenin equivalents are not limited to single source circuits. It is possible to find the equivalent of a network with several sources. Finding the open circuit output voltage will undoubtedly require some extra work, for example the use of superposition of source conversions. Finding the Thevenin resistance is unchanged, just remember to replace every source with its ideal internal resistance before simplifying the network.
            \bigbreak \noindent 
            Consider 
            \bigbreak \noindent 
            \fig{.8}{./figures/42.png}
            \bigbreak \noindent 
            Suppose we want to determine the Thevenin equivalent of the circuit driving the 1k$\Omega$. Then,
            \bigbreak \noindent 
            \fig{.8}{./figures/43.png}
            \bigbreak \noindent 
            Now, since there is no current passing through the $200$, we have
            \begin{align*}
                R_{\text{EQ}} = 100 + 800 = 900
            .\end{align*}
            Thus,
            \begin{align*}
                I = \frac{10}{900} = 0.0111A
            .\end{align*}
            Since there is no voltage drop across the 200, $E_{th}$ is given by the voltage at $b$, which is equal to the voltage drop across the 800, or the source voltage minus the voltage drop across the 100.
            \begin{align*}
                10 - V_{100} &= 10 - 0.0111(100) = 8.89V, \\
                0.0111(800) &= 8.89V
            .\end{align*}
            The equivalent resistance is found by shorting the voltage source and then simplifying the circuit. We have:
            \begin{align*}
                200 + 100 || 800 = 288.89 \Omega
            .\end{align*}
            Now, we can find $V_{c}$
            \bigbreak \noindent 
            \fig{.4}{./figures/44.png}
            \bigbreak \noindent 
            In this equivalent circuit,
            \begin{align*}
                I = \frac{8.889}{1000 + 288.89} = 0.006897A
            .\end{align*}
            Thus,
            \begin{align*}
                V_{c} = V_{1k} = 0.006987(1000) = 6.897V
            .\end{align*}
            Now, lets analyse the original circuit. We have 
            \begin{align*}
                R_{\text{EQ}} = (1000 + 200) || 800 + 100 = 580 \Omega
            .\end{align*}
            Thus,
            \begin{align*}
                I &= \frac{10}{580} = 0.0172A, \\
                V_{b} &= 10 - 0.0172(100) = 8.276V
            .\end{align*}
            The current through $c$ is given by 
            \begin{align*}
                I_{c} = I_{200} = I_{1000} = \frac{V_{b}}{1200} = 0.0069A
            .\end{align*}
            Thus,
            \begin{align*}
                V_{c} = 0.0069(1000) = 6.897V
            .\end{align*}
            Notice that the voltages found at $c$ match perfectly, while the voltage at $b$ is only slightly off.
        \item \textbf{Norton's theorem}: In a nutshell, Norton's theorem is the current source version of Thevenin's theorem. 
            \bigbreak \noindent 
            That is, a single port DC network can be reduced to a single current source, $I_{N}$, with parallel internal resistance, $R_{N}$.
            \bigbreak \noindent 
            replace all sources with their ideal internal resistance and then perform appropriate series and parallel combinations to reduce this to a single resistance value. Consequently, the Norton and Thevenin resistance values are identical, $R_{N} = R_{th} $. Second, instead of finding the open circuit output voltage, we find the short circuit output current. This is the Norton current. Instead of thinking in terms of a voltmeter at the opened load, we think in terms of connecting a shorting ammeter across the load. Either way, we're looking for the maximum value that can be obtained.
            \bigbreak \noindent 
            Perhaps the most useful thing to remember here is that if we can create a Thevenin equivalent for a network then it must be possible to create a Norton equivalent. Indeed, once a Thevenin equivalent is found, a source conversion can be performed on it to yield the Norton equivalent!
        \item \textbf{Normalized source}: Consider a source voltage $E$ with internal resistance $R_{i}$ connected in series to load resistor $R$. Suppose we wanted to normalize such that $E = R_{i} = 1$.
            \bigbreak \noindent 
            Define
            \begin{align*}
                E_{\text{base}} = E, \quad R_{\text{base}} = R_{i}
            .\end{align*}
            Define normalized quantities by dividing each value by its corresponding base
            \begin{align*}
                e = \frac{E}{E_{\text{base}}} = 1, \quad r_{i} = \frac{R_{i}}{R_{\text{base}}} = 1
            .\end{align*}
            The normalized load resistance becomes 
            \begin{align*}
                r = \frac{R}{R_{i}}
            .\end{align*}
            Normalization does not physically change \(E\) or \(R_i\); it expresses the circuit using \(E\) volts and \(R_i\) ohms as the chosen units.
        \item \textbf{Base}: A base is a chosen reference value used to express a physical quantity as a dimensionless ratio. 
            \bigbreak \noindent 
            For any quantity $X$,
            \begin{align*}
                X_{\text{normalized}} = \frac{X_{\text{actual}}}{X_{\text{base}}}
            .\end{align*}
            To recover an actual quantity, multiply its normalized value by its base
        \item \textbf{Maximum Power Transfer Theorem}: Given a simple voltage source with internal resistance, a useful question to ask is “What value of load resistance will yield the maximum amount of power in the load?” While it is not true that maximizing load power is a goal of all circuit designs, it is a goal of a portion of them and thus worth a closer look.
            \bigbreak \noindent 
            \fig{.6}{./figures/45.png}
            \bigbreak \noindent 
            We have source $E$, source internal resistance $R_{i}$ and load resistance $R$. 
            \bigbreak \noindent 
            We would like to describe the load power in terms of the load resistance. To make the job easier, we may normalize the voltage source $E$ to 1 volt and the source resistance $R_i$ to 1 Ohm. By doing this, $R$ also becomes a normalized value, that is, it no longer represents a simple resistance value but rather represents a ratio in comparison to $R_i$. In this way the analysis will work for any set of source values. Note that the value of $E$ will equally scale the power in both $R_i$ and $R$, so a precise value is not needed, and thus, we may as well chose 1 volt for convenience.
            \bigbreak \noindent 
            Note that 
            \begin{align*}
                P = I^{2}R, \quad I = \frac{E}{R_{i} + R} = \frac{1}{1 +  R}
            .\end{align*}
            Hence,
            \begin{align*}
                P = \left(\frac{1}{1+r}\right)^{2}R = \frac{R}{R^{2} + 2R + 1}
            .\end{align*}
            We now have an equation that describes the load power in terms of the load resistance. Notice that if $R = 0$ or $R = \infty$ (shorted and open case) the power is zero.
            \bigbreak \noindent 
            While shorting the load will yield the maximum load current, it will also yield zero load voltage and thus no power. Similarly, opening the load will yield maximum load voltage but it will also yield zero load current, and again, no load power
            \bigbreak \noindent 
            Its easy to see that maximum power will be achieved when $R = 1$. Thus,
            \begin{quote}
               "\textit{Maximum load power will be achieved when the load resistance is equal to the internal resistance of the driving source. }" 
            \end{quote}
            While matching the resistance produces the maximum load power, it does not produce maximum load current or maximum load voltage. In fact, this condition produces a load voltage and a load current that are half of their maximums. Their product, however, is at the maximum. Further, efficiency at maximum load power is only 50\% (i.e., only half of all generated power goes to the load with the other half being wasted internally). Values of $R$ greater than $R_{i}$ will achieve higher efficiency but at reduced load power. Sometimes we favor efficiency over maximal load power.
        \item \textbf{Delta-Y Conversions}: Certain component configurations, such as bridged networks, cannot be reduced to a single resistance using basic series-parallel conversion techniques. One method for simplification involves converting sections into more convenient forms. The configurations in question are three-point networks containing three resistors. Due to the manner in which they drawn, they are referred to as delta networks and $Y$ networks
            \bigbreak \noindent 
            \fig{.6}{./figures/46.png}
            \bigbreak \noindent 
            Alternately, if they are slightly redrawn they are known as pi networks and $T$ networks
            \bigbreak \noindent 
            \fig{.6}{./figures/47.png}
            \bigbreak \noindent 
            It is possible to convert back and forth between delta and $Y$ networks. That is, for every delta network, there exists a $Y$ network such that the resistances seen between the $X$, $Y$ and $Z$ terminals are identical, and vice versa. Consequently, one configuration can replace another in order to simplify a larger circuit.
        \item \textbf{$\Delta$-Y conversions}: A true equivalent circuit would present the same resistance between any two terminals as the original circuit. Consider the unloaded case. The equivalent resistances seen between each pair of terminals for the delta and $Y$ respectively are
            \begin{align*}
                R_{XY} &= R_{a} || (R_{b} + R_{c}) = R_{d} + R_{e} \\
                R_{XZ} &= R_{b} || (R_{a} + R_{c}) = R_{d} + R_{f} \\
                R_{ZY} &= R_{c} || (R_{b} + R_{a}) = R_{e} + R_{f} 
            .\end{align*}
            Assuming we have the delta and are looking for the $Y$ equivalent, note that we have three equations with three unknowns ($R_{d}$, $R_{e}$ and $R_{f}$). Thus, they can be solved using a term elimination process.
            \bigbreak \noindent 
            If we subtract the third equation from the first, we eliminate the second resistance ($R_{e}$) and arrive at a difference between the first and third unknown resistances ($R_{d} - R_{f}$). This quantity can then be added to the second to eliminate the third resistance ($R_{f}$), leaving just the first unknown resistance ($R_{d}$).
            \begin{align*}
                (R_d + R_e) - (R_e + R_f) &= R_d - R_f  = \left(R_a \parallel (R_b + R_c)\right) - \left(R_b \parallel (R_a + R_c)\right), \\[6pt]
                (R_d + R_f) + (R_d - R_f) &= 2R_d = R_b \parallel (R_a + R_c) + R_a \parallel (R_b + R_c) - R_c \parallel (R_a + R_b) 
            .\end{align*}
            Hence,
            \begin{align*}
                2R_{d} = R_{b} || (R_{a} + R_{c}) + R_{a} || (R_{b} + R_{c}) - R_{c} || (R_{a} + R_{b})
            .\end{align*}
            After simplifying,
            \begin{align*}
                R_{d} = \frac{R_{a}R_{b}}{R_{a} + R_{b} + R_{c}}
            .\end{align*}
            Similarly, 
            \begin{align*}
                R_{e} &= \frac{R_{a}R_{c}}{R_{a} + R_{b} + R_{c}}, \\
                R_{f} &= \frac{R_{b}R_{c}}{R_{a} + R_{b} + R_{c}}
            .\end{align*}
        \item \textbf{$Y-\Delta$ conversion}: First, observe that
            \begin{align*}
                \frac{R_{d}}{R_{e}} = \frac{R_{b}}{R_{c}}
            .\end{align*}
            Thus,
            \begin{align*}
                R_{b} = \frac{R_{c}R_{d}}{R_{e}}
            .\end{align*}
            Similarly, 
            \begin{align*}
                R_{a} &= \frac{R_{d}R_{c}}{R_{f}}, \\
                R_{c} &= \frac{R_{a}R_{f}}{R_{d}}
            .\end{align*}
            Now, we can plug the expressions for $R_{a}$ and $R_{b}$ into 
            \begin{align*}
                R_{d} = \frac{R_{a}R_{b}}{R_{a} + R_{b} + R_{c}}
            .\end{align*}
            Thus,
            \begin{align*}
                R_{d} = \frac{\frac{R_{d}^{2}R_{c}^{2}}{R_{f}R_{e}}}{\frac{R_{d}R_{c}}{R_{f}} + \frac{R_{d}R_{c}}{R_{e}} + R_{c}}
            .\end{align*}
            Yielding
            \begin{align*}
                R_{c} = \frac{R_{d}R_{e} + R_{d}R_{f} + R_{e}R_{f}}{R_{e}}
            .\end{align*}
            Now, substitute this into
            \begin{align*}
                R_{b} = \frac{R_{c}R_{d}}{R_{e}}
            \end{align*}
            and
            \begin{align*}
                R_{a} = \frac{R_{c}R_{d}}{R_{f}}
            .\end{align*}
            Hence,
            \begin{align*}
                R_{a} &= \frac{R_{d}R_{e} + R_{e}R_{f} + R_{d}R_{f}}{R_{f}} \\
                R_{b} &= \frac{R_{d}R_{e} + R_{e}R_{f} + R_{d}R_{f}}{R_{e}} \\
                R_{c} &= \frac{R_{d}R_{e} + R_{e}R_{f} + R_{d}R_{f}}{R_{d}}
            .\end{align*}
            If three identical resistors are used, the values of the delta equivalent will all be three times that value, the inverse of the situation when converting from delta to Y


    \end{itemize}

    \pagebreak 
    \unsect{Nodal and Mesh Analysis, Dependent Sources}
    \begin{itemize}
        \item \textbf{Dependent sources}: Current and voltage sources whos value is not fixed to some particular value, but rather is dependent on some other current or voltage in the circuit.
        \item \textbf{Nodal and mesh analysis}: Nodal analysis and mesh analysis are two systematic methods used to solve for unknown voltages and currents in DC circuits.
        \item \textbf{Nodal analysis}: Nodal analysis is a technique that can be applied to virtually any circuit. In general use, it might be considered a universal solution technique as there are no practical circuit configurations that it cannot handle. Nodal analysis relies on the application of Kirchhoff's current law to create a series of node equations that can be solved for node voltages. These equations are based on Ohm's law and will be of the form $I = V/R$
            \bigbreak \noindent 
            Once the node voltages are obtained, finding any branch currents or component powers becomes an almost trivial exercise. We will examine two variations; a general version that can be used with both voltage and current sources, and a second somewhat quicker version that can be used with circuits only driven by current sources.
        \item \textbf{Nodal analysis, general method}: Consider
            \bigbreak \noindent 
            \fig{.6}{./figures/49.png}
            \bigbreak \noindent 
            We begin by labeling connection nodes and assigning current directions. We are particularly interested in \textbf{current junctions}, that is, places where currents can combine or split. These are also known as \textbf{summing nodes} and are circled in green on the figure. The current directions are chosen arbitrarily and may be the opposite of reality. This is not a problem. If we assign directions that are wrong, the resulting directions will ultimately show up reversed but the computed node voltages will be just fine.
            \bigbreak \noindent 
            One node is chosen as the reference. This is the point to which all other node voltages will be measured against. Typically, the reference node is ground, although it does not have to be.
            \bigbreak \noindent 
            We now write a current summation equation for each summing node, except for the reference node. In this circuit there is only one node where currents combine (other than ground) and that's node $b$. Points $a$ and $c$ are places where components connect, but they are not summing nodes, so we can ignore them for now. Using KCL on node $b$ we can say:
            \begin{align*}
                I_{1} + I_{2} = I_{3}
            .\end{align*}
            Next, we describe these currents in terms of the node voltages and associated components via Ohm's law. For example, $I_{3}$ is the node $b$ voltage divided by $R_{3}$ while $I_{1}$ is the voltage across $R_{1}$ divided by $R_{1}$. This voltage is $V_{a} - V_{b}$. Therefore
            \begin{align*}
                \frac{V_{a}-V_{b}}{R_{1}} + \frac{V_{c} - V_{b}}{R_{2}} = \frac{V_{b}}{R_{3}}
            .\end{align*}
            With $V_{a} = E_{1}$, and $V_{c} = E_{2}$ and some algebra, we have
            \begin{align*}
               \left(\frac{1}{R_{1}}\right) E_{1} + \left(\frac{1}{R_{2}}\right)E_{2} = \left(\frac{1}{R_{1}} + \frac{1}{R_{2}} + \frac{1}{R_{3}}\right)V_{b}
            .\end{align*}
            Notice that we have a series of products of conductances and voltages. All quantities are known except for $V_{b}$ and thus it is easily found with a little more algebra.
            \bigbreak \noindent 
            For current sources, a more direct approach is possible. We start as before, identifying nodes and labeling currents.
            \bigbreak \noindent 
            \fig{.6}{./figures/50.png}
            \bigbreak \noindent 
            We then write current summation equations at each node (except for ground). We consider currents entering a node as positive, and exiting as negative.
            \begin{align*}
                \text{Node a: } I1 &= I3 + I4 \\
                \text{Node b: } I3 &= I2 + I5, \text{ and rearranging in terms of the fixed source,} \\
                \text{Node b: } -I2 &= -I3 + I5
            .\end{align*}
            The currents are then described by their Ohm's law equivalents:
            \begin{align*}
                \text{Node a: } I_{1} &= \frac{V_{a} - V_{b}}{R_{3}} + \frac{V_{a}}{R_{1}} \\
                \text{Node b: } -I_{2} &= -\frac{V_{a} - V_{b}}{R_{3}} + \frac{V_{b}}{R_{2}}
            .\end{align*}
            Hence,
            \begin{align*}
                \text{Node a: } I_{1} &= \left(\frac{1}{R_{1}} + \frac{1}{R_{3}}\right)V_{a} - \left(\frac{1}{R_{3}}\right)v_{b} \\
                \text{Node b: } -I_{2} &= -\left(\frac{1}{R_{3}}\right)V_{a} + \left(\frac{1}{R_{3}} + \frac{1}{R_{2}}\right)V_{b}
            .\end{align*}
            As the resistor values and currents are known, simultaneous equation solution techniques may be used to solve for the node voltages. Once again, there will be as many equations as node voltages
        \item \textbf{Inspection Method}: The system of equations can be obtained directly through inspection if the circuit contains no voltage sources
            \bigbreak \noindent 
             For the node under inspection, sum all of the current sources connected to it to obtain the current constant The conductance term for that node will be the sum of all of the conductances connected to that node. For the other node conductances, determine the conductances between the node under inspection and these other nodes. These terms will all be negative As a crosscheck, the set of equations produced must exhibit diagonal symmetry
             \bigbreak \noindent 
             The inspection method is summarized as follows:
             \begin{enumerate}
                 \item Verify that the circuit uses only current sources with resistors and no voltage sources. If voltage sources exist, they must be converted to current sources before proceeding
                 \item Find all of the current summing nodes and number them. Also decide on the reference node (usually ground).
                 \item To generate an equation, locate the first node. This is the node of interest and the next few steps will be associated with it.
                 \item Sum the current sources feeding the node of interest. Entering is deemed positive while exiting is deemed negative. The sum is placed on one side of the equals sign
                 \item Next, find all of the resistors connected to the node of interest and write them as a sum of conductances on the other side of the equals sign, the group being multiplied by this node's voltage (e.g., $V_{1}$).
                 \item Now find all of the resistors that are connected to the node of interest and to other nodes (except the ground reference). For each of these other nodes, multiply the sum of the conductances between the node of interest and this other node by this other node's voltage, and then subtract that product from the equation built so far. Once all other nodes are considered, this equation is finished.
                 \item Find the next node and treat this as the new node of interest
                 \item Repeat steps 4 through 7 until all nodes have been treated as the node of interest. Each iteration creates a new equation. There will be as many equations as there are nodes, less the reference node. Check for diagonal symmetry and solve.
             \end{enumerate}
             As an example, consider
             \bigbreak \noindent 
             \fig{.6}{./figures/51.png}
             \bigbreak \noindent 
             Observe the three nodes $a$, $b$ and reference ground. Focus on node $a$. The current sources feeding this node are 800mA in, 2mA out. Thus,
             \begin{align*}
                 0.8A - 2A = ...
             .\end{align*}
             The conductances connected to node $a$ give
             \begin{align*}
                 0.8A - 2A = \left(\frac{1}{10} + \frac{1}{50} + \frac{1}{100}\right) V_{a}...
             .\end{align*}
             Now find all of the resistors that are connected to this node and to the other nodes. Multiply those resistors (expressed as conductances) by the other associated node voltages and subtract the products from the expression built so far. Repeat for all remaining nodes except the ground reference. In this example there is only one other node, node $b$, and thus only one iteration.
             \begin{align*}
                 0.8A - 2A &= \left(\frac{1}{10} + \frac{1}{50} + \frac{1}{100}\right) V_{a} - \left(\frac{1}{50} + \frac{1}{100}\right) V_{b} \\
                  &= -1.2A = 130mSV_{a} - 30mSV_{b}
             .\end{align*}
             Recall that conductance is expressed in \textbf{siemens} (S).
             \bigbreak \noindent 
             Now we repeat the entire process for the next equation. Node $b$ is our new node of interest. The fixed current sources are:
             \begin{align*}
                 -800mA - 300mA &= -\left(\frac{1}{50} + \frac{1}{100}\right)V_{a} + \left(\frac{1}{25} + \frac{1}{50} + \frac{1}{100}\right)V_{b} \\
                                &= -1.1A = -30mSV_{a} + 70mSV_{b}
             .\end{align*}
             Hence, our two equations are 
             \begin{align*}
                 -1.2A &= 130mSV_{a} - 30mSV_{b} \\
                 -1.1A &= -30mSV_{a} + 70mSV_{b}
             .\end{align*}
             So,
             \begin{align*}
                 \begin{pmatrix} 0.130 & -0.030 \\ -0.030 & 0.070 \end{pmatrix}\begin{pmatrix} V_{a} \\ V_{b} \end{pmatrix} = \begin{pmatrix} -1.2 \\ -1.1 \end{pmatrix}
             .\end{align*}
             Notice that we indeed have symmetry. Solving this system yields
             \begin{align*}
                 V_{a} &= -14.27V, \\
                 V_{b} &= -21.83V
             .\end{align*}
             The current flowing through the 100 $\Omega$ resistor is 
             \begin{align*}
                 \frac{V_{a} - V_{b}}{100} = 75.6mA
             \end{align*}
             flowing left to right.
             \bigbreak \noindent 
             Next, consider
             \bigbreak \noindent 
             \fig{.6}{./figures/52.png}
             \bigbreak \noindent 
             We get
             \begin{align*}
                 \begin{pmatrix} 0.4  & -0.25 & 0 \\ -0.25 & 0.8 & -0.5 \\ 0 & -0.5 & 0.6 \end{pmatrix} \begin{pmatrix} V_{a} \\ V_{b} \\ V_{c} \end{pmatrix} = \begin{pmatrix} 1.5 \\ 0 \\ 1 \end{pmatrix}
             .\end{align*}
             Thus,
             \begin{align*}
                 V_{a}  &= 8.62V \\
                 V_{b} &= 7.8V \\
                 V_{c} &= 8.17V
             .\end{align*}
            \item \textbf{Converting Sources and Other Simplifications}: Given circuits with voltage sources, it may be easier to convert them to current sources and then apply the inspection technique rather than using the general approach outlined initially. There is one trap to watch out for when using source conversions: the voltage across or current through a converted component will most likely not be the same as the voltage or current in the original circuit. This is because the location of the converted component will have changed. 
                \bigbreak \noindent 
                Consider
                \bigbreak \noindent 
                \fig{.6}{./figures/53.png}
                \bigbreak \noindent 
                could be solved using nodal analysis by converting the voltage source and the associated resistance into a current source. That is, $E/R_{1}$ would be converted into a source $I_{3}$ with a parallel resistor $R_{1}$.
                \bigbreak \noindent 
                \fig{.6}{./figures/54.png}
                \bigbreak \noindent 
                In the converted circuit, although $R_1$ still connects to node $a$, the other end no longer connects to the voltage source. Rather, the right side now connects to node $c$. Therefore, the voltage drop across $R_1$ in the converted circuit is not likely to equal the voltage drop seen across $R_1$ in the original circuit (unless $E$ was 0 volts).
                \bigbreak \noindent 
                In the converted circuit, nodes $a$ and $c$ have not changed from the original, so the original voltage across $R_1$ can be determined via Va, $V_{c}$ and $E$ in the original circuit
                \bigbreak \noindent 
                If a circuit uses voltage sources exclusively, or even a very large proportion of voltage sources, an alternate technique called mesh analysis may be a better choice than using numerous conversions. 
            \item \textbf{Reduction}: A final item to consider is to simplify resistor networks in order to reduce the number of nodes. Fewer nodes means fewer equations and a quicker solution.
            \item \textbf{Supernode}: Consider
                \bigbreak \noindent 
                \fig{.6}{./figures/55.png}
                \bigbreak \noindent 
                On occasion you may come across a circuit like this, a voltage source without a series resistance. Without that resistance, it becomes impossible to create an expression for the current passing through the source using the general method, and impossible to convert the voltage source into a current source in order to use the inspection method.
                \bigbreak \noindent 
                One possible way out of this quandary is to simply add a very small resistor in series with it so that a source conversion is possible. The resistor in question would have to be much smaller than any surrounding resistors in order to have minimal impact on the results. A reduction by two orders of magnitude would generally yield a variation smaller than that produced by resistor tolerances in all but high precision circuits. The other way out is to use a \textbf{supernode}.
                \bigbreak \noindent 
                A supernode is, in effect, the combination of two nodes. It relies on a simple observation. Notice that the path of the voltage source produces identical currents flowing into and out of nodes $a$ and $b$
                \bigbreak \noindent 
                As a consequence, if we treat the two nodes as one big node, then when we write a KCL summation, these two terms will cancel.
                \bigbreak \noindent 
                \textbf{Note:} For a two-terminal voltage source, the current leaving the positive terminal has the same magnitude as the current entering the negative terminal. The current may split into several paths throughout the external circuit, but all those branch currents eventually recombine before returning to the negative terminal. Charge cannot continuously accumulate inside the source or circuit.
                \bigbreak \noindent 
                This describes a voltage source \textbf{delivering power}. If current instead enters the positive terminal and leaves the negative terminal, the source is \textbf{absorbing power}, as when a battery is being charged.
                \bigbreak \noindent 
                \fig{.6}{./figures/56.png}
                \bigbreak \noindent 
                In this version we have replaced the voltage source with its ideal internal resistance, a short. We have also labeled the two nodes of interest, $a$ and $b$, and labeled the currents, drawn with convenient directions
                \bigbreak \noindent 
                Due to the shorted voltage source, nodes $a$ and $b$ are now the same node. Consider the currents entering and exiting this combined or “super” node. On the left side (formerly node $a$) we see a constant current Ix entering while $I_{1}$, $I_{2}$ and $I_{3}$ are exiting. On the right side (formerly node $b$) we see $I_{y}$ exiting along with $I_{4}$, and entering we see $I_{1}$ and $I_{2}$. Now let's set up the entering currents on the left side of the equals sign with the exiting currents on the right:
                \begin{align*}
                    \sum I_{\text{in}} = I_{x} + I_{1} + I_{2} = \sum I_{\text{out}} = I_{y} + I_{1} + I_{2} + I_{3} + I_{4}
                .\end{align*}
                Thus,
                \begin{align*}
                    I_{x} - I_{y} = I_{3} + I_{4} = \frac{1}{R_{1}}V_{a} + \frac{1}{R_{3}}V_{b}
                .\end{align*}
                We also know that $V_{a} - V_{b} = E$ from the original circuit. Assuming all sources and resistors are known, that makes two equations with two unknowns, solvable using simultaneous equation techniques. 
                \bigbreak \noindent 
                Consider
                \bigbreak \noindent 
                \fig{.6}{./figures/57.png}
                \bigbreak \noindent 
                The KCL equation for node $a$ is 
                \begin{align*}
                    1 + I_{E} - \frac{V_{a} - V_{b}}{20} - \frac{V_{a}}{4} = 0
                .\end{align*}
                The KCL equation for node $b$ is 
                \begin{align*}
                    -2.5 + \frac{V_{a}  - V_{b}}{20} - I_{E} - \frac{V_{b}}{10} = 0
                .\end{align*}
                Adding these two yields
                \begin{align*}
                    -\frac{V_{a}}{4}-1.5-\frac{V_{b}}{10} = 0
                .\end{align*}
                Hence,
                \begin{align*}
                    \frac{V_{a}}{4} + \frac{V_{b}}{10} = -1.5
                .\end{align*}
                The second equation is $V_{a} - V_{b} = 60$. Thus,
                \begin{align*}
                    \begin{pmatrix} 0.25 & 0.1 \\ 1 & -1  \end{pmatrix}\begin{pmatrix} V_{a} \\ V_{b} \end{pmatrix} = \begin{pmatrix} -1.5 \\ 60 \end{pmatrix}
                .\end{align*}
                Therefore,
                \begin{align*}
                    V_{a} &= 12.857V \\
                    V_{b} &= -47.143V
                .\end{align*}
            \item \textbf{Planar circuit}: A planar circuit is an electrical circuit that can be drawn on a flat, two-dimensional surface without any of its wires or branches crossing each other. The defining feature is that a crossing-free drawing can exist
                \bigbreak \noindent 
                In essence, these circuits can be drawn so that they appear like a series of window panes.
            \item \textbf{Branch current and mesh current}: Branch currents are the actual, physical currents flowing through individual components in a circuit, while mesh currents are mathematical loop currents used to simplify circuit analysis
                \bigbreak \noindent 
                A branch current is a hypothetical circulating current assigned to a closed loop or "window" of a planar circuit. It is purely conceptual; shared elements carry a net combination of adjacent mesh currents rather than a single standalone mesh current
            \item \textbf{Mesh analysis}: In some respects mesh analysis is a mirror of nodal analysis. While nodal analysis leverages KCL to create a series of node equations that are used to solve for node voltages, mesh analysis uses KVL to create a series of loop equations that can be solved for mesh currents. 
                \bigbreak \noindent 
                While branch currents represent the current flowing through a particular component, mesh currents combine to create branch currents. That is, the current through any particular component may be an individual mesh current or a combination of two mesh currents.
                \bigbreak \noindent 
                Circuits using complex series-parallel arrangement with multiple voltage and/or current sources may be solved using this technique. It is, however, limited to planar circuits
                \bigbreak \noindent 
                Non-planar circuits may be 3D in appearance and in such a circuit it becomes impossible to define the loop equations as any given component might have more than two meshing currents.
            \item \textbf{General method}:  We begin by designating a series of loops. These loops should be minimal in size and cover all components at least once. By convention, the loops are drawn clockwise. There is nothing magic about them being clockwise, it is just a matter of consistency. The current path along a loop is referred to as a mesh current.
                \bigbreak \noindent 
                \fig{.6}{./figures/59.png}
                \bigbreak \noindent 
                we have two loops, and thus, two mesh currents, $I_1$ and $I_2$. Note that all components exist in at least one loop (and sometimes in more than one loop, like $R_3$). Depending on circuit values, one or more of these loop's directions may in fact be opposite of reality. This is not a problem. If this is the case, the currents will show up as negative values, and thus we know that they're really flowing counter-clockwise.
                \begin{align*}
                    \text{Loop 1: } E_{1} &= \text{ voltage across } R_{1} + \text{ voltage across } R_{3} \\
                    \text{Loop 2: } -E_{2} &= \text{ voltage across } R_{2} + \text{ voltage across } R_{3}
                .\end{align*}
                $E_2$ is negative as $I_2$ is drawn flowing out of its negative terminal. If we choose downward as the positive reference direction, then \(I_1\) is positive and \(I_2\) is negative. Using Ohm's law:
                \begin{align*}
                    E_{1} &= I_{1}R_{1} + (I_{1} - I_{2})R_{3} \\
                    -E_{2} &= I_{2}R_{2} + (I_{2} - I_{1})R_{3}
                .\end{align*}
                For the second loop, we take upward current to be positive.
                \begin{itemize}
                    \item For loop 1, you traverse \(R_3\) downward, so the current in that direction is
                    \[
                        I_{R_3,\downarrow}=I_1-I_2.
                    \]
                \item For loop 2, you traverse \(R_3\) upward, so the current in that direction is
                    \[
                        I_{R_3,\uparrow}=I_2-I_1.
                    \]
            \end{itemize}
            These are the same physical current expressed using opposite reference directions:
            \begin{align*}
                I_{R_3,\uparrow}=-I_{R_3,\downarrow}
            .\end{align*}
            Therefore, it is valid to use downward as positive while writing loop 1 and upward as positive while writing loop 2. You simply must remain consistent within each equation.
            \bigbreak \noindent 
            We have
            \begin{align*}
                E_{1} &= (R_{1} + R_{3})I_{1} - R_{3}I_{2} \\
                -E_{2} &= R_{3}I_{1} + (R_{2}+R_{3})I_{2}
            .\end{align*}
            Assuming that the resistor values and source voltages are known, we have two equations with two unknowns. These can be solved for $I_1$ and $I_2$ using simultaneous equation solution techniques such as determinants or Gauss-Jordan elimination.
            \bigbreak \noindent 
            Consider an example
            \bigbreak \noindent 
            \fig{.8}{./figures/60.png}
            \bigbreak \noindent 
            We have
            \begin{align*}
                \ell_{1}:\; 12V = (I_{1} - I_{2})1k + (I_{1}-I_{3})2k \\
                \ell_{2}:\; 0V = (I_{2}-I_{1})1k + 3kI_{2}+(I_{2} - I_{3})5k \\
                \ell_{3}:\; 0V = (I_{3} - I_{2})5k + 4kI_{3} + (I_{3} - I_{1})2k
            .\end{align*}
            Thus, the system is 
            \begin{align*}
                1000\begin{pmatrix} 3 & -1 & -2 \\-1 & 9 & -5 \\ -2 & -5 & 11 \end{pmatrix}\begin{pmatrix} I_{1} \\ I_{2} \\ I_{3} \end{pmatrix} = \begin{pmatrix} 12 \\ 0 \\ 0 \end{pmatrix}
            .\end{align*}
            Hence,
            \begin{align*}
                I_{1} = 5.729mA \\
                I_{2} = 1.626mA \\
                I_{3} = 1.781mA 
            .\end{align*}
            To find $V_{b}$, we use net current through the $2k$ resistor. Thus,
            \begin{align*}
                V_{b} = (I_{1} - I_{3})2000 = 7.8968V
            .\end{align*}
            Similarly, 
            \begin{align*}
                V_{bc} = (I_{3} - I_{2})5000 = 774mV
            .\end{align*}
        \item \textbf{Inspection Method}: Like nodal analysis, there is a method to generate the equations by inspection. Simply focus on one loop, which we'll call the loop under inspection, and ask the following questions: what is the total source voltage in this loop? This yields the voltage constant for the equation. Then sum all of the resistance values in the loop under inspection. This yields the coefficient for that current term. For the remaining current coefficients in this equation, sum the resistances that are in common between the loop under inspection and the other loops, respectively. These values will always be negative. Repeat this process for all loops, with each loop in turn becoming the loop under inspection. As was the case with nodal analysis, the set of equations produced must exhibit diagonal symmetry
            \bigbreak \noindent 
            While it is possible to extend this technique to include current sources, it is often easier and less error-prone to convert the current sources into voltage sources and continue with the direct inspection method outlined above. Either way, it is important to remember that the number of loops determines the number of equations to be solved.
            \bigbreak \noindent 
            Consider
            \bigbreak \noindent 
            \fig{.6}{./figures/61.png}
            \bigbreak \noindent 
            We have
            \begin{align*}
                50V &= (100 + 200)I_{1} - 200I_{2} = 300I_{1} - 200I_{2} \\
                0V &= (300 + 400 + 200)I_{2} - 200I_{1} - 400I_{3} = - 200I_{1} + 900I_{2}  - 400I_{3} \\
                20V &= (500 + 400)I_{3} - 400I_{2} = 0I_{1} - 400I_{2} + 900 I_{3}
            .\end{align*}
            Thus,
            \begin{align*}
                \begin{pmatrix} 300 & -200 & 0 \\ -200 & 900 & -400 \\ 0 & -400 & 900 \end{pmatrix}\begin{pmatrix} I_{1} \\ I_{2} \\ I_{3} \end{pmatrix} = \begin{pmatrix} 50 \\ 0 \\ 20 \end{pmatrix}
            .\end{align*}
        \item \textbf{Supermesh}: Sometimes you may run across a current source which has no associated internal resistance
            \bigbreak \noindent 
            This is similar to the situation seen previously under nodal analysis where a voltage source does not have a specified internal resistance. There are two ways of solving this predicament. The first is add a very large resistance in parallel with the current source and then perform a source conversion on the pair so that the inspection method of mesh can be followed. The larger the value of this resistor, the more accuracy that is obtained. As a general rule it should be, at minimum, at least a couple of orders of magnitude larger than any surrounding resistor, and preferably larger. The second method is to use supermesh. A supermesh is a larger mesh loop that contains other mesh loops inside of it.
            \bigbreak \noindent 
            \fig{.6}{./figures/62.png}
            \bigbreak \noindent 
            In the center we have a current source, Is, which lacks an associated internal resistance. Two traditional mesh loops, $I_1$ and $I_2$, are labeled as usual. The problem here is that we cannot use an Ohm's law based $IR$ voltage drop for Vb. We have no way to express this as the voltage across Is is an unknown. On the other hand, what we do know is that Is must equal the combination of the traditional mesh currents $I_1$ and $I_2$. That is, from the perspective of the first loop, $I_{s} = I_{1} - I_{2}$. Remember, one or both of the mesh currents could be negative, and thus rotating counterclockwise.
            \bigbreak \noindent 
            At this point we invoke the idea of a supermesh loop. First, we replace the offending current source with its ideal internal resistance (an open). The supermesh loop is drawn which encompasses the original two loops.
            \bigbreak \noindent 
            \fig{.6}{./figures/63.png}
            \bigbreak \noindent 
            We now perform a KVL summation around the supermesh loop, similar to what we have done previously. The difference this time is that we need to recognize that components each see one of the original mesh currents; namely $I_1$ or $I_2$ in this case. We do not solve for a supermesh current, we simply use the supermesh to define the loop for the KVL summation
            \begin{align*}
                \sum V_{\text{rises}} = E_{1} = \sum V_{\text{drops}} = V_{R_{1}} + V_{R_{2}} + E_{2}
            .\end{align*}
            The voltage drops across the resistors can be expanded using Ohm's law, using the original mesh current associated with each resistor.
            \begin{align*}
                E_{1} - E_{2} = I_{1}R_{1} + I_{2}R_{2}
            .\end{align*}
            By inspection,
            \begin{align*}
                I_{S} = I_{1} - I_{2} 
            .\end{align*}
            We now have two equations with two unknowns and can solve for $I_1$ and $I_2$.
            \bigbreak \noindent 
            Consider
            \bigbreak \noindent 
            \fig{.6}{./figures/64.png}
            \bigbreak \noindent 
            So, we have
            \bigbreak \noindent 
            \fig{.6}{./figures/65.png}
            \bigbreak \noindent 
            By KVL, we have
            \begin{align*}
                20 = 5I_{1} + 8I_{2} + 12
            .\end{align*}
            So, 
            \begin{align*}
                8 = 5I_{1} + 8I_{2}
            .\end{align*}
            By inspection, 
            \begin{align*}
                3 = -I_{1} - I_{2}
            .\end{align*}
            Thus,
            \begin{align*}
                \begin{pmatrix} 5 & 8 \\ -1 & 1 \end{pmatrix}\begin{pmatrix} I_{1} \\ I_{2} \end{pmatrix} = \begin{pmatrix} 8 \\ 3 \end{pmatrix}
            .\end{align*}
            So,
            \begin{align*}
                I_{1} &= -1.231A \\
                I_{2} &= 1.769A
            .\end{align*}
        \item \textbf{Dependent Sources}: A dependent source is a current or voltage source whose value is not fixed (i.e., independent) but rather which depends on some other circuit current or voltage. The general form for the value of a dependent source is
            \begin{align*}
                Y = kX
            \end{align*}
            where $X$ and $Y$ are currents or voltages and $k$ is a proportionality factor. For example, the value of a dependent voltage source may be a function of a current, so instead of the source being equal to, say, 10 volts, it could be equal to twenty times the current passing through a particular resistor, or $V=20I$.
            \bigbreak \noindent 
            There are four possible dependent sources: the voltage-controlled voltage source (VCVS), the current-controlled voltage source (CCVS), the voltage-controlled current source (VCCS), and the current-controlled current source (CCCS).
            \bigbreak \noindent 
            The source and control parameters are the same for both the VCVS and the CCCS so k is unitless (although it may be given as volts/volt and amps/amp, respectively). For the VCCS and CCVS, k has units of amps/volt and volts/amp, respectively. These are referred to as the \textbf{transconductance} and \textbf{transresistance} of the sources with units of siemens and ohms, respectively.
            \bigbreak \noindent 
            The schematic symbols for dependent or controlled sources are usually drawn using a diamond. Also, there may be a secondary connection for the controlling current or voltage.
            \bigbreak \noindent 
            \fig{.6}{./figures/66.png}
            \bigbreak \noindent 
            On each of these symbols, the control element is shown to the left of the source. This portion is not always drawn on a schematic. Instead, the source simply may be labeled as a function. For example,
            \begin{align*}
                V = 0.02 I_{x}
            \end{align*}
            where $I_{x}$i s the controlling current. Dependent sources are not “off-the-shelf” items in the same way that a battery is. Rather, dependent sources usually are used to model the behavior of more complex devices. For example, a bipolar junction transistor commonly is modeled as a CCCS while a field effect transistor may be modeled as a VCCS.
            \bigbreak \noindent 
            Similarly, many op amp circuits are modeled as VCVS systems. Solutions for circuits using dependent sources follow along the lines of those established for independent sources (i.e., the application of Ohm's law, KVL, KCL, etc.), however, the sources are now dependent on the remainder of the circuit which tends to complicate the analysis.
        \item \textbf{Isolated and coupled}: In general, there are two possible configurations: isolated and coupled. A example of the isolated form is
            \bigbreak \noindent 
            \fig{.6}{./figures/67.png}
            \bigbreak \noindent 
            In this configuration, the dependent source (center) does not interact with the subcircuit on the left driven by the independent source, thus it can be analyzed as two separate circuits. Solutions for this form are relatively straightforward in that the control value for the dependent source can be computed directly. This value is then substituted into the dependent source and the analysis continues as is usual. Sometimes it is convenient if the solution for a particular voltage or current is defined in terms of the control parameter rather than as a specific value (e.g., the voltage across a particular resistor might be expressed as 8 $V_{A}$ instead of 12 volts, where $V_{A}$ is 1.5 volts).
            \bigbreak \noindent 
            The second type of circuit (coupled) is somewhat more complex in that the dependent source can affect the parameter that controls the dependent source. In other words, the dependent source(s) will contribute terms that include the controlling parameter(s), thus partly controlling itself. Some additional effort will be in order to solve these circuits.
            \bigbreak \noindent 
            \fig{.6}{./figures/68.png}
            \bigbreak \noindent 
            In this example it should be obvious that the current from the dependent source can affect the voltage at node a, and it is this very voltage that in turn sets up the value of the current source. Circuits of this type can be analyzed using mesh or nodal analysis. Nodal analysis works well here and is illustrated following.
            \bigbreak \noindent 
            We begin by defining current directions. Assume that the currents through $R_1$ and $R_3$ are flowing into node $a$, the current through $R_2$ is flowing out of node $a$, and the current through $R_4$ is flowing out of node $b$. We shall number the branch currents to reflect the associated resistor.
            \begin{align*}
                &\text{Node $a$: } I_{1} + I_{3} = I_{2} \\
                &\text{Node $b$: } kV_{a} = I_{3} + I_{4}
            .\end{align*}
            Notice that 
            \begin{align*}
                I_{1} &= \frac{E - V_{a}}{R_{1}} \\
                I_{2} &= \frac{V_{a}}{R_{2}} \\
                I_{3} &= \frac{V_{b} - V_{a}}{R_{3}} \\
                I_{4} &= \frac{V_{b}}{R_{4}}
            .\end{align*}
            Thus,
            \begin{align*}
                &\text{Node $a$: } \frac{E - V_{a}}{R_{1}} + \frac{V_{b}-V_{a}}{R_{3}} = \frac{V_{a}}{R_{2}} \\
                &\text{Node $b$: } kV_{a} = \frac{V_{b} - V_{a}}{R_{3}} + \frac{V_{b}}{R_{4}}
            .\end{align*}
            Hence,
            \begin{align*}
                &\text{Node $a$: } \frac{E}{R_{1}} - \frac{V_{a}}{R_{1}} + \frac{V_{b}}{R_{3}} - \frac{V_{a}}{R_{3}} = \frac{V_{a}}{R_{2}} \\
                &\text{Node $b$: } kV_{a} = \frac{V_{b}}{R_{3}} - \frac{V_{a}}{R_{3}} + \frac{V_{b}}{R_{4}}
            .\end{align*}
            So,
            \begin{align*}
                &\text{Node $a$: } \frac{E}{R_{1}} = \left(\frac{1}{R_{1}} + \frac{1}{R_{2}} + \frac{1}{R_{3}}\right)V_{a} - \frac{1}{R_{3}}V_{b} \\
                &\text{Node $b$: } 0 = -\left(k + \frac{1}{R_{3}}\right)V_{a} + \left(\frac{1}{R_{3}} + \frac{1}{R_{4}}\right)V_{b}
            .\end{align*}
            Values for the resistors, $k$, and $E$ normally are known, so the analysis proceeds directly.
            \bigbreak \noindent 
            Also, it is worth remembering that it is possible to perform source conversions on dependent sources, within limits. The same procedure is followed as for independent sources. The new source will remain a dependent source (e.g., converting VCVS to VCCS). This process is not applicable if the control parameter directly involves the internal impedance (i.e., is its voltage or current).
            \bigbreak \noindent 
            Consider
            \bigbreak \noindent 
            \fig{.8}{./figures/69.png}
            \bigbreak \noindent 
            First observe that $a$ and $b$ refer to the same node. This CCCS would be typical of a simple model of a bipolar junction transistor (nodes $a$, $b$, and $c$). Ideally, the output voltage, $V_{c}$, would be equal to the input voltage (1 V) times the resistor ratio of 15 k$\Omega$ / 2 k$\Omega$ and inverted, or approximately -7.5 volts. In reality, it usually comes up a little short of this. Let's see how well this works out.
            \bigbreak \noindent 
            Note that
            \begin{align*}
                I_{x} = \frac{1V - V_{b}}{10k}
            .\end{align*}
            Using KCL at node $b$, we have
            \begin{align*}
                I_{x} + 100I_{x} = 101I_{x} = 101 \frac{1-V_{b}}{10000} = \frac{V_{b}}{2000}
            .\end{align*}
            Thus,
            \begin{align*}
                \frac{101}{10k} = \left(\frac{1}{2k} + \frac{101}{10k}\right)V_{b}
            .\end{align*}
            Hence,
            \begin{align*}
                V_{b} = 0.95283V
            .\end{align*}
            For node $c$:
            \begin{align*}
                -\frac{V_{c}}{15k} = 100I_{x} = 100 \frac{1-V_{b}}{10k}
            .\end{align*}
            So,
            \begin{align*}
                V_{c} = -7.0755V
            .\end{align*}


    \end{itemize}

    \pagebreak 
    \unsect{Capacitors}
    \begin{itemize}
        \item \textbf{Intro}: Capacitors are fundamentally different from resistors in terms of both their construction and their operation. For starters, when placed in DC circuits, capacitors are not ohmic, unlike resistors. Their current-voltage characteristic does not respond to Ohm's law. Instead, their current-voltage characteristic is dynamic in nature. Further, in the ideal case, capacitors do not dissipate power. Indeed, capacitors are energy storage devices. In a way, you could imagine them to be a little like rechargeable batteries, but unlike batteries these devices could be charged or discharged in tiny fractions of a second and last for billions upon billions of cycles.         
            \bigbreak \noindent 
            Capacitors are an integral part of modern electronic systems. They are used in AC-to-DC power supplies to help smooth and stabilize the output voltage. In audio and communications systems they are used in filters, for example to control the high and low frequency response of amplifiers and similar equipment, or for tuning purposes. Essentially, in any application that needs to smooth out a varying voltage, store electric charge or filter a signal; a capacitor is likely to be used.
        \item \textbf{Capacitance and Capacitors}: A capacitor is a device that stores energy. Capacitors store energy in the form of an electric field. At its most simple, a capacitor can be little more than a pair of metal plates separated by air. As this constitutes an open circuit, DC current will not flow through a capacitor. If this simple device is connected to a DC voltage source. Negative charge will build up on the bottom plate while positive charge builds up on the top plate. This process will continue until the voltage across the capacitor is equal to that of the voltage source. In the process, a certain amount of electric charge will have accumulated on the plates.
            \bigbreak \noindent 
            \fig{.8}{./figures/70.png}
            \bigbreak \noindent 
            The ability of this device to store charge with regard to the voltage appearing across it is called \textbf{capacitance}. Its symbol is $C$ and it has units of \textbf{farads (F)}.
            \bigbreak \noindent 
            By definition, if a total charge of 1 coulomb is associated with a potential of 1 volt across the plates, then the capacitance is 1 farad
            \begin{align*}
                1 \text{ farad } \equiv 1 \text{ coulomb } / 1 \text{ volt}
            .\end{align*}
            That is,
            \begin{align*}
                C = \frac{Q}{V}
            .\end{align*}
            For any given voltage, the greater the capacitance, the greater the amount of charge that can be stored. We can also see that, given a certain size capacitor, the greater the voltage, the greater the charge that is stored. These observations relate directly to the amount of energy that can be stored in a capacitor.
            \bigbreak \noindent 
            Unsurprisingly, the energy stored in capacitor is proportional to the capacitance. It is also proportional to the square of the voltage across the capacitor.
            \begin{align*}
                W = \frac{1}{2}CV^{2}
            .\end{align*}
            The basic capacitor consists of two conducting plates separated by an insulator, or dielectric. This material can be air or made from a variety of different materials such as plastics and ceramics.
            \bigbreak \noindent 
            \fig{.8}{./figures/71.png}
            \bigbreak \noindent 
            For practical capacitors, the plates may be stacked alternately or even made of foil and formed into a rolled tube. However it is constructed, the characteristics of the dielectric will play a major role in the performance of the device, as we shall see.
            \bigbreak \noindent 
            In general, capacitance increases directly with plate area, $A$, and inversely with plate separation distance, $d$. Further, it is also proportional to a physical characteristic of the dielectric; the \textbf{permittivity}, $\varepsilon$. Thus, capacitance is equal to:
            \begin{align*}
                C = \varepsilon \frac{A}{d}
            .\end{align*}
        \item \textbf{Dielectric}: A dielectric is an electrical insulator that can store electrical energy when placed in an electric field through a process called polarization. 
        \item \textbf{Polarization}: When exposed to an external electric field, the positive and negative charges inside the material shift slightly. Positive charges move toward the field, and negative charges move in the opposite direction, creating tiny internal electric dipoles. This is \textbf{polarization}.
            \bigbreak \noindent 
            In polarization, the charges cannot escape because they are tightly bound to the atoms or molecules of the dielectric material. Unlike conductors, dielectrics do not have "free electrons."
        \item \textbf{Dipoles}: A dipole is a separation of positive and negative charges (or magnetic poles) over a small distance, creating an overall neutral object or molecule that has distinct positive and negative ends.
        \item \textbf{Permittivity}: Dielectric permittivity measures how much a material polarizes and stores electrical energy when placed in an external electric field
        \item \textbf{Fringing}: It should be noted that the effective plate area is somewhat larger than the precise physical area of the plates. This is due to a phenomenon called \textbf{fringing}. Essentially, the electric field lines bulge outward at the plate edges rather than maintain uniform parallel orientation.
            \bigbreak \noindent 
            \fig{.8}{./figures/72.png}
        \item \textbf{Efficiency}: the permittivity of the dielectric plays a major role in determining the \textbf{volumetric efficiency} of the capacitor, in other words, the amount of capacitance that can be packed into a given sized component. Some dielectrics are notably more efficient than others. To make comparisons easier, \textbf{relative permittivity} is often used, that is, the ratio of the dielectric's permittivity to that of a vacuum, $\varepsilon_{0}$
            \bigbreak \noindent 
            It might seem that choosing the dielectric with the highest permittivity would be the best choice but this is not necessarily the case. There are several other factors that go into this decision including temperature stability, leakage resistance (effective parallel resistance), ESR (equivalent series resistance) and breakdown strength. For an ideal capacitor, leakage resistance would be infinite and ESR would be zero.
        \item \textbf{Breakdown strength}: Unlike resistors, capacitors do not have maximum power dissipation ratings. Instead, they have maximum voltage ratings. The \textbf{breakdown strength} of the dielectric will set an upper limit on how large of a voltage may be placed across a capacitor before it is damaged.
            \bigbreak \noindent 
            Breakdown strength is measured in volts per unit distance, thus, the closer the plates, the less voltage the capacitor can withstand. For example, halving the plate distance doubles the capacitance but also halves its voltage rating.
        \item \textbf{Capacitor Styles and Packaging}: Unlike resistors, whose physical size relates to their power rating and not their resistance value, the physical size of a capacitor is related to both its capacitance and its voltage rating
            \bigbreak \noindent 
            Modest surface mount capacitors can be quite small while the power supply filter capacitors commonly used in consumer electronics devices such as an audio amplifier can be considerably larger than a D cell battery
            \bigbreak \noindent 
            For large capacitors, the capacitance value and voltage rating are usually printed directly on the case. Some capacitors use “MFD” which stands for “microfarads”. While a capacitor color code exists, rather like the resistor color code, it has generally fallen out of favor. For smaller capacitors a numeric code is used that echoes the color code. Typically it consists of a three digit number such as “152”.
            \bigbreak \noindent 
            The first two digits are the precision portion and the third digit is the power of ten multiplier. The result is in picofarads. Thus, 152 is 1500 pf.
            \bigbreak \noindent 
            \fig{.6}{./figures/73.png}
            \bigbreak \noindent
            \fig{.6}{./figures/74.png}
        \item \textbf{Capacitor Data Sheet}: 
            \bigbreak \noindent 
            \fig{.6}{./figures/75.png}
        \item \textbf{Capacitors in Series and in Parallel}: Multiple capacitors placed in series and/or parallel do not behave in the same manner as resistors. Placing capacitors in parallel increases overall plate area, and thus increases capacitance,
            \bigbreak \noindent 
            Therefore capacitors in parallel add in value, behaving like resistors in series. In contrast, when capacitors are placed in series, it is as if the plate distance has increased, thus decreasing capacitance. Therefore capacitors in series behave like resistors in parallel. Their value is found via the reciprocal of summed reciprocals or the product-sum rule.
            \bigbreak \noindent 
            If a circuit contains nothing but a voltage source in parallel with a group of capacitors, the voltage will be the same across all of the capacitors, just as it is in a resistive parallel circuit. If the circuit instead consists of multiple capacitors that are in series with a voltage source
            \bigbreak \noindent 
            the voltage will divide between them in inverse proportion. In other words, the larger the capacitance, the smaller its share of the applied voltage. The voltages can also be found by first determining the series equivalent capacitance. The total charge may then be determined using the applied voltage. 
            \bigbreak \noindent 
            Finally, the individual voltages are computed from
            \begin{align*}
                V = \frac{Q}{C}
            .\end{align*}
            Consider
            \bigbreak \noindent 
            \fig{.6}{./figures/76.png}
            \bigbreak \noindent 
            We have
            \begin{align*}
                C_{\text{total}} &= \frac{1}{\frac{1}{C_{1}} + \frac{1}{C_{2}} + \frac{1}{C_{3}}} =1.143\mu F
            .\end{align*}
            Hence,
            \begin{align*}
                Q = VC = 12V (1.143)\mu F = 13.71\mu C
            .\end{align*}
            Charge is constant across all of the series capacitors, so
            \begin{align*}
                V_{2\mu F} &= \frac{13.71}{2} = 6.855V, \\
                V_{4\mu F} &= \frac{13.71}{4} = 3.427V, \\
                V_{8\mu F} &= \frac{13.71}{8} = 1.714V
            .\end{align*}
            The sum of the three voltages is 12 volts (within rounding error) and verifies KVL as expected.
        \item \textbf{Note:} While it may be tempting to try, do not attempt to verify the operation of Example 8.3 in the laboratory using a standard DMM. The reason is because the internal resistance of a typical digital voltmeter is many orders of magnitude lower than the leakage resistance of the capacitors. As a result, charge will be transferred to the meter, ruining the measurement. It would be akin to trying to measure the voltages across a string of resistors, each in excess of 100 M$\Omega$, with a meter whose internal resistance is 1 M$\Omega$. The meter's resistance dominates the parallel combination and causes excessive loading which ruins the measurement. A special sort of voltmeter, an electrostatic voltmeter or electrometer, is needed for these types of measurements. These are sometimes referred to as non-charge transfer meters.
        \item \textbf{Current-Voltage Relationship}: The fundamental current-voltage relationship of a capacitor is not the same as that of resistors. Capacitors do not so much resist current; it is more productive to think in terms of them reacting to it. The current through a capacitor is equal to the capacitance times the rate of change of the capacitor voltage with respect to time (i.e., its slope). That is, the value of the voltage is not important, but rather how quickly the voltage is changing. Given a fixed voltage, the capacitor current is zero and thus the capacitor behaves like an open. If the voltage is changing rapidly, the current will be high and the capacitor behaves more like a short. Expressed as a formula
            \begin{align*}
                i = C \frac{dv}{dt}
            .\end{align*}
            Note that $i$ is the current flowing through the capacitor. From this, we see 
            \begin{align*}
                \frac{dv}{dt} = \frac{i}{C}
            .\end{align*}
            If a capacitor is fed by a constant current source, the voltage will rise at a constant rate ($dv/dt$). It is continuously depositing charge on the plates of the capacitor at a rate of $I$, which is equivalent to $Q/t$. As long as the current is present, feeding the capacitor, the voltage across the capacitor will continue to rise. A good analogy is if we had a pipe pouring water into a tank, with the tank's level continuing to rise. This process of depositing charge on the plates is referred to as charging the capacitor.
            \bigbreak \noindent 
            \fig{.6}{./figures/77.png}
            \bigbreak \noindent 
            \fig{.6}{./figures/78.png}
            \bigbreak \noindent 
            As time progresses, the voltage across the capacitor increases with a positive polarity from top to bottom. With a theoretically perfect capacitor and source, this would continue forever, or until the current source was turned off. In reality, this line would either begin to deflect horizontally as the source reached its limits, or the capacitor would fail once its breakdown voltage was reached. The slope of this line is dictated by the size of the current source and the capacitance
            \bigbreak \noindent 
            Consider
            \bigbreak \noindent 
            \fig{.6}{./figures/79.png}
            \bigbreak \noindent 
            We wish to determine the rate of change of voltage across the capacitor, and determine the capacitors voltage 10 ms after power is switched on.
            \begin{align*}
                \frac{dv}{dt} = \frac{i}{C} = \frac{-5\mu A}{30nF} = -166.7 \frac{V}{s}
            .\end{align*}
            So, after 10 ms, we see
            \begin{align*}
                -166.7(10) = -1.667V
            .\end{align*}
            This is assuming it is completely uncharged when power is applied.
        \item \textbf{Limit}: If a capacitor is driven by a fixed current source, the voltage across it rises at the constant rate of $i/C$. There is a limit to how quickly the voltage across the capacitor can change. An instantaneous change means that $dv/dt$ is infinite, and thus, the current driving the capacitor would also have to be infinite (an impossibility). This is not an issue with resistors, which obey Ohm's law, but it is a limitation of capacitors. Therefore we can state a particularly important characteristic of capacitors:
            \begin{quote}
               "\textit{The voltage across a capacitor cannot change instantaneously.}" 
            \end{quote}
            Note that instantaneous change means the voltage changes in zero time. Thus, the rate becomes infinitely large. This would force the current $i$ to also be infinite.
            \bigbreak \noindent 
            This observation will be key to understanding the operation of capacitors in DC circuits.
        \item \textbf{Initial and Steady-State Analysis of RC Circuits}: When analyzing resistor-capacitor circuits, always remember that capacitor voltage cannot change instantaneously. If we assume that a capacitor in a circuit is not initially charged, then its voltage must be zero. 
            \bigbreak \noindent 
            The instant the circuit is energized, the capacitor voltage must still be zero. If there is no voltage across the device, then it is behaving like a short circuit. We call this the initial state. Thus, we have our first rule regarding RC circuits:
            \begin{quote}
                "\textit{For DC analysis, initially capacitors appear as shorts}"
            \end{quote}
            Assume that $C_{1}$ and $C_{2}$ are initially uncharged and there is no voltage across them.
            \bigbreak \noindent 
            \fig{.6}{./figures/80.png}
            \bigbreak \noindent 
            The instant power is applied, the two capacitors appear as short circuits. If we redraw the circuit for this instant in time
            \bigbreak \noindent 
            \fig{.6}{./figures/81.png}
            \bigbreak \noindent 
            Given this equivalent, we can see that shorting $C_2$ places $R_2$ and $R_3$ in parallel, however, they are both shorted out by $C_1$. This leaves only $R_1$ left in the circuit along with the source, $E$. At this point, currents will begin to flow, and thus begin charging up the capacitors. As the capacitor voltages rise, the current will begin to decrease, and eventually the capacitors will stop charging. At that point no further current will be flowing, and thus the capacitor will behave like an open. We call this the \textbf{steady-state condition} and we can state our second rule:
            \begin{quote}
               "\textit{At steady-state, capacitors appear as opens.}" 
            \end{quote}
            Current decreases as a capacitor charges because the voltage building up across the capacitor's plates grows closer to the supply voltage, reducing the effective voltage difference that pushes electrons through the circuit
            \bigbreak \noindent 
            The current flowing into a capacitor is directly proportional to how fast its voltage changes over time (\(i = C \frac{dv}{dt}\)). As the capacitor approaches full charge, the voltage changes more slowly (\[\frac{dv}{dt}\] approaches zero), which forces the current down to zero.
            \bigbreak \noindent 
            at steady-state both capacitors behave as opens
            \bigbreak \noindent 
            \fig{.8}{./figures/82.png}
            \bigbreak \noindent 
            This leaves $E$ to drop across $R_1$ and $R_2$. This will create a simple voltage divider. The steady-state voltage across $C_1$ will equal that of $R_2$. As $C_2$ is also open, the voltage across $R_3$ will be zero while the voltage across $C_2$ will be the same as that across $R_2$.
            \bigbreak \noindent 
            In reality, practical capacitors can be thought of as an ideal capacitance in parallel with a very large (leakage) resistance, so there will be a limit to this performance.
            \bigbreak \noindent 
            Consider
            \bigbreak \noindent 
            \fig{.8}{./figures/83.png}
            \bigbreak \noindent 
            Initially, the capacitor is treated as a short. Thus, 
            \begin{align*}
                V_{6k} = 24 - \frac{24}{6k||3k+1k}(1000) = 16V
            .\end{align*}
            When the capacitor is fully charged, we treat it as an open. Thus,
            \begin{align*}
                V_{6k}= 24 - \frac{24}{1k + 6k}(1000) = 20.57V
            .\end{align*}
            \textbf{Remark}: A capacitor acts like a short circuit only at the exact moment \(t=0\) when it is completely uncharged, and transitions gradually into an open circuit as the voltage across it builds up
        \item \textbf{Time constant resistance}: Use the equivalent resistance seen by the capacitor, i.e the Thevenin resistance
            \begin{align*}
                \tau = R_{\text{th}}C
            .\end{align*}
        \item \textbf{Note on the time constant $\tau$}: Recall that 
            \begin{align*}
                C = \frac{Q}{V}
            ,\end{align*}
            with
            \begin{align*}
                Q = It
            .\end{align*}
            Thus,
            \begin{align*}
                C = \frac{It}{V}
            .\end{align*}
            Now, with $V = IR$, 
            \begin{align*}
                \frac{1}{I} = \frac{R}{V}
            .\end{align*}
            Thus,
            \begin{align*}
                \tau = RC = \frac{R}{V}It = \frac{1}{I}It = t
            .\end{align*}
            So, $\tau$ is indeed a measure of time.
        \item \textbf{Transient Response of RC Circuits}: The question remains, “What happens between the time the circuit is powered up and when it reaches steady-state?” This is known as the transient response
            \bigbreak \noindent 
            \fig{.8}{./figures/84.png}
            \bigbreak \noindent 
            The key to the analysis is to remember that capacitor voltage cannot change instantaneously. Assuming the capacitor is uncharged, the instant power is applied, the capacitor voltage must be zero. Therefore all of the source voltage drops across the resistor. This creates the initial current, and this current starts to charge the capacitor
            \bigbreak \noindent 
            According to Kirchhoff's voltage law, as the capacitor voltage begins to increase, the resistor voltage must decrease because the sum of the two must equal the fixed source voltage. This means that the circulating current must also decrease
            \bigbreak \noindent 
            This, in turn, means that the rate of capacitor voltage increase begins to slow. As the capacitor voltage continues to increase, less voltage is available for the resistor, causing further reductions in current, and a further slowing of the rate of capacitor voltage change. Eventually, the capacitor voltage will be nearly equal to the source voltage.
            \bigbreak \noindent 
            This will result in a very small potential across the resistor and an equally small current, slowing subsequent capacitor voltage increases to a near standstill.
            \bigbreak \noindent 
            Theoretically, the capacitor voltage approaches the source voltage but never quite equals it.
            \bigbreak \noindent 
            Similarly, the current drops to near zero, but never completely turns off.
            \bigbreak \noindent 
            \fig{.4}{./figures/85.png}
            \bigbreak \noindent 
            The dashed red line represents the initial rate of change of capacitor voltage. This trajectory is what would be expected if an ideal current source drove the capacitor. As noted previously, the rate of voltage change versus times is equal to $i/C$, and therefore in this case, $E/(RC)$. If the initial rate of change were to continue unabated, the source voltage, $E$, would be reached in $RC$ seconds. Consequently, $RC$ is referred to as the charge \textbf{time constant} and is denoted by $\tau$. Thus
            \begin{align*}
                \tau = RC
            .\end{align*}
            As noted, once the capacitor begins to charge, the current begins to decrease and the capacitor voltage curve begins to fall away from the initial trajectory. The solid red curve represents the capacitor voltage. Notice that after five time constants the capacitor is nearly fully charged and the circuit is considered to be in steady-state
            \begin{quote}
               "\textit{Steady-state is reached in approximately five time constants.}" 
            \end{quote}
            This sort of recursively dependent operation is characteristic of exponential functions. The equation for the capacitor's voltage charging curve is:
            \begin{align*}
                V_{C}(t) = E\left(1-e^{-\frac{t}{\tau}}\right)
            .\end{align*}
            The dashed blue line shows the initial slope of current change. The solid blue curve shows the circulating current (and by extension of Ohm's law, the resistor voltage). The equation for this curve, which follows the general shape $e^{-t}$ is 
            \begin{align*}
                I(t) &= \frac{E}{R}e^{-\frac{t}{\tau}}, \\
                V_{R}(t) &= Ee^{-\frac{t}{\tau}}
            .\end{align*}
            Once power is removed or bypassed, the stored charge on the capacitor will dissipate through any associated resistor(s) creating a discharge current which will end with the capacitor voltage drained back to zero. During the discharge phase, both the capacitor's voltage and current will follow the solid blue curve
            \bigbreak \noindent 
            The discharge time constant may be different from the charge times constant, depending on the associated resistances.
            \bigbreak \noindent 
            Consider
            \bigbreak \noindent 
            \fig{.8}{./figures/86.png}
            \bigbreak \noindent 
            Assume the switch is closed at time $t = 0$. Determine the charging time constant, the amount of time after the switch is closed before the circuit reaches steady-state, and the capacitor voltage at $t = 0$, $t = 50$ milliseconds and $t = 1$ second. Assume the capacitor is initially uncharged.
            \bigbreak \noindent 
            The time constant is 
            \begin{align*}
                \tau = RC = 50k\Omega \cdot 2\mu F = 100ms
            .\end{align*}
            Steady-state will be reached in five time constants, or 500 milliseconds. So,
            \begin{align*}
                V_{C}(0) &= 0V,\\
                V_{C}(1) &= 100V
            .\end{align*}
            Also,
            \begin{align*}
                V_{C}(50) = 100V\left(1-e^{-\frac{50}{100}}\right) \approx 39.35V
            .\end{align*}
            \bigbreak \noindent 
            Next, consider
            \bigbreak \noindent 
            \fig{.6}{./figures/87.png}
            \bigbreak \noindent 
            At time $t = 0$ the switch contacts position 1. The switch is thrown to position 2 at time $t = 50$ milliseconds.
            \bigbreak \noindent 
            Determine the charging time constant, the amount of time after the switch is closed before the circuit reaches steady-state, the maximum charging and discharging currents, and the capacitor voltage at $t = 0$, $t = 50$ milliseconds, $t = 90$ milliseconds, and $t = 1$ second. 
            \bigbreak \noindent 
            We have 
            \begin{align*}
               \tau = 20k\Omega(220nF) = 4.4ms
            .\end{align*}
            Thus, steady-state will be reached in 
            \begin{align*}
                5\tau = 5(4.4) = 22ms
            .\end{align*}
            Thus,
            \begin{align*}
                V_{C}(0) &= 0,\\
                V_{C}(50) &= 12
            .\end{align*}
            The maximum charging current will occur at $t=0$ when all of the 12 volt source drops across the 20k$\Omega$ resistor, or 600$\mu A$, flowing left to right.
            \bigbreak \noindent 
            At $t=50$, the switch is thrown to position two. The 12 volt source and 20 k$\Omega$ resistor are no longer engaged. At this point the capacitor has 12 volts across it, positive to negative, top to bottom. As the capacitor voltage cannot change instantaneously, the capacitor now acts as a voltage source and discharges through the 120 k$\Omega$ resistor
            \bigbreak \noindent 
            Note that the discharge current is flowing counterclockwise, the opposite of the charging current. The discharge time constant is:
            \begin{align*}
                \tau_{\text{discharge}} = RC = 120000 \Omega \cdot 220 \times 10^{-9}F = 26.4ms
            .\end{align*}
            The capacitor will fully discharge down to 0 volts in 5 time constants, or some 132 milliseconds after the switch is thrown to position 2.
            \bigbreak \noindent 
            Thus steadystate occurs at $t = 182$ milliseconds. The maximum discharge current occurs the instant the switch is thrown to position 2 when all of the capacitor's 12 volts drops across the 120 k$\Omega$ resistor, yielding 100 $\mu A$, flowing top to bottom.
            \bigbreak \noindent 
            Clearly, at $t = 90$ milliseconds the capacitor is in the discharge phase. The capacitor's voltage and current during the discharge phase follow the solid blue curve of Figure 8.25. The elapsed time for discharge is 90 milliseconds minus 50 milliseconds, or 40 milliseconds net.
            \begin{align*}
                V_{C}(t) = Ee^{-\frac{t}{\tau}} = 12 - e^{-\frac{40}{26.4}} = 2.637V
            .\end{align*}
            \bigbreak \noindent 
            Basic single resistor-capacitor circuits prove to be fairly easy to solve given a little practice, but what if a more complex circuit is used? In this situation the section feeding the capacitor may be simplified using Thevenin's theorem to determine the effective source voltage and charging resistance. The circuit then reverts back to a simple RC network which may be solved directly
            \bigbreak \noindent 
            If power is interrupted before the capacitor is fully charged, the equations presented previously may be used to determine the precise voltage(s) and current(s) reached. The capacitor will then behave as a voltage source and begin to discharge.
            \bigbreak \noindent 
            Consider
            \bigbreak \noindent 
            \fig{.8}{./figures/98.png}
            \bigbreak \noindent 
            Determine the charging time constant, the amount of time after the switch is closed before the circuit reaches steady-state, and the capacitor voltage at $t = 0$, 100 milliseconds, and 200 milliseconds. At 200 milliseconds, the switch is opened. Determine how long it takes for the capacitor to fully discharge and the voltage across the 6 k$\Omega$ resistor at $t = 275$ milliseconds (i.e., 75 milliseconds after the switch is opened).
            \bigbreak \noindent 
            Let's use Thevenin. If we replace the capacitor with an open, we see that the Thevenin voltage $E_{th}$ is the voltage dropped by $R_{6k}$. With source current
            \begin{align*}
                I_{s} = \frac{24}{7000} = 0.003428A
            ,\end{align*}
            we have
            \begin{align*}
                E_{\text{th}} = 0.003428(6000) = 20.57V
            .\end{align*}
            Now, the resistance $R_{\text{th}}$ is 
            \begin{align*}
                R_{\text{th}} = 6k || 1k + 3k = 3.857
            .\end{align*}
            Thus,
            \begin{align*}
                \tau = 3.857 \times 10^{3} \cdot 10 \times 10^{-6} = 38.57ms
            .\end{align*}
            So,
            \begin{align*}
               5\tau = 192.85ms 
            .\end{align*}
            Thus, the capacitor will reach 20.57V in 192.85ms.
            \begin{align*}
                V_{C}(0) = 0V, \quad V_{C}(200) = 20.57V
            .\end{align*}
            At 100ms,
            \begin{align*}
                V_{C}(100) = E(1-e^{-\frac{t}{\tau}}) = 20.57\left(1-e^{-\frac{100}{38.57}}\right) = 19.03V
            .\end{align*}
            With the voltage source removed, the capacitor discharges.
            \begin{align*}
                \tau_{\text{discharge}} = RC = 9k\Omega(10\mu F) = 90ms
            .\end{align*}
            Steady-state after discharge will be reached 
            \begin{align*}
                5\tau_{\text{discharge}} = 450ms
            .\end{align*}
            Since the switch is flipped open at $t=200$, this occurs at $t=650ms$. Regarding $V_{6k}$ at $ t=275ms$,
            \begin{align*}
                V_{C}(275 - 200) = V_{C}(75) = 20.57\left(1-e^{-\frac{75}{90}}\right) = 8.94V
            .\end{align*}
            The current through the 6k and the 3k is 
            \begin{align*}
                I = \frac{8.94}{6000 + 3000}
            .\end{align*}
            Thus,
            \begin{align*}
                V_{6k} = \frac{8.94}{9000}\cdot 6000 = 5.96V
            .\end{align*}


    \end{itemize}

    \pagebreak 
    \unsect{Inductors}
    \begin{itemize}
        \item \textbf{Intro}: Inductors are fundamentally different from both resistors and capacitors in terms of their construction and their operation. Inductors do, however, share certain broad traits with capacitors. First, they are energy storage devices. In the case of the inductor, energy is stored in a magnetic field, similar to the case of the capacitor which utilizes an electric field. Further, in the ideal case, inductors do not dissipate power. Also like capacitors, when inductors are placed in DC circuits they are not ohmic, meaning that their current-voltage characteristic does not respond to Ohm's law. Instead, their currentvoltage characteristic is dynamic in nature. In some respects, though, their current-voltage behavior is opposite to the way in which capacitors behave, and thus they offer their own unique performance characteristics.
            \bigbreak \noindent 
            Unfortunately, real-world inductors generally do not behave as close to their desired ideal operation as do realworld capacitors. The secondary effects of inductor construction limit their performance; perhaps the most notable factor being their potentially large equivalent series resistance. They are also susceptible to external magnetic fields which can introduce noise and interference, degrading signal quality. For these reasons, there are areas where, given a choice, capacitors would be preferred over inductors. But this is by no means a broad condemnation and there are areas where the use of inductors is essential. Beyond this, the very concept of inductance is important in that informs designers of the practical limits of performance of their circuits.
            \bigbreak \noindent 
            Inductors are commonly referred to as coils or chokes, for reasons that will be apparent shortly.
        \item \textbf{Unit and symbol}: The unit of inductance is the henry ($H$), the symbol for inductance is $L$.
        \item \textbf{Inductance and Inductors}: To begin, we need to examine the interrelation between electric current and magnetic fields in a conductor. When a current passes through a conductor, such as a wire, a magnetic field is created around the conductor that is proportional to the strength of the current.
            \bigbreak \noindent 
            \fig{.8}{./figures/99.png}
            \bigbreak \noindent 
            The magnetic field can be thought of as sets of concentric rings around the conductor, although for clarity only single loops are drawn. The number of magnetic lines in a given area is known as the \textbf{magnetic flux} and is given the symbol $\Phi$. The unit of megnetic flux is the \textbf{weber} (wb).
            \begin{align*}
                \text{Magnetic flux } \equiv \text{ the number of magnetic lines enclosed in a given area}
            .\end{align*}
            Note that the magnetic field runs the length of the conductor. The direction of the field lines follows the right hand rule: if you grasp the wire with your right hand such that your thumb is pointing in the direction of conventional current flow, then your fingers wrap in the direction of the magnetic field.
        \item \textbf{Electromagnetic induction}: A changing magnetic field induces an electric current in a conductor, a process known as electromagnetic induction.
            \bigbreak \noindent 
            A steady, unmoving magnetic field produces no current at all. The magnetic field must change over time to create electricity. This change involves the amount of magnetic field passing through a closed loop of wire, known as magnetic flux
        \item \textbf{Solenoid}: If we form the conductor into a loop, the field lines are corralled into the center of the loop.
            \bigbreak \noindent 
            \fig{.6}{./figures/100.png}
            \bigbreak \noindent 
            The enhancement effect can be magnified by adding more loops in tandem. This is known as a solenoid It is the most basic form of an inductor.
            \bigbreak \noindent 
            A solenoid is an electromechanical device that converts electrical energy into linear mechanical motion using a coil of wire and a movable metal core
            \bigbreak \noindent 
            An electric current flows through a tightly wound coil of copper wire. The current creates a uniform, directional magnetic field inside the coil. The magnetic field pulls a movable metal core (called a plunger or armature) into the center of the coil. A spring pushes the plunger back to its original resting position when the electric current is turned off.
            \bigbreak \noindent 
            \bigbreak \noindent 
            \fig{.8}{./figures/102.png}
            \bigbreak \noindent 
            \fig{.8}{./figures/101.png}
            \bigbreak \noindent 
            In this figure, the coil is shown from the side, as a cross-section of the individual loops. The dots inside the conductors indicate that the current is flowing toward you, out of the page; while the the crosses indicate that the current is flowing into the page. The lines of flux exit out the right, loop around, and reenter from the left. Due to limited space, the entire loop for each line is not drawn, and it is important to remember that magnetic flux lines do not end, but always create a loop. Further, although it is shown as a plane here, this field is three dimensional, with lines looping back into the page as well as in front of it.
            \bigbreak \noindent 
            This is how electromagnets can be created. The north pole is the exiting end (right side), while the south pole is the entering end (left side). 
            \bigbreak \noindent 
            If the current changes, there will be a commensurate change in the magnetic field. Further, this change in the field will induce a current in the conductor that creates a magnetic field that opposes the original change in the field This is known as Lenz's law. Alternately, it can be stated that the induced current caused by a changing magnetic field will oppose the change in the original current that created that change in the original magnetic field.
            \bigbreak \noindent 
            At this point we can offer a proper definition of the weber:
            \begin{quote}
               \textit{1 weber $ \equiv $ the magnetic flux that, acting on a single loop of a conductor, produces a potential of 1 volt if the flux is reduced to zero at an even rate over 1 sec}
            \end{quote}
        \item \textbf{Induced current}: An induced current is an electric current produced in a conductor when it experiences a changing magnetic environment
        \item \textbf{Magnetic flux density}: In magnetic circuits we are also interested in the magnetic flux density which is the magnetic flux per unit area. The symbol for flux density is $B$ and has units of teslas ($T$),
            \begin{align*}
                1 \text{ tesla } \equiv 1 \text{ weber } / \text{ meter }^{2}
            .\end{align*}
        \item \textbf{Inductance}: Finally, we come to the definition of inductance and its unit, the henry.
            \bigbreak \noindent 
            Inductance is a measure of the tendency of a conductor to oppose a change in the current flowing through it.
            \begin{align*}
                1 \text{ henry } \equiv  1 \text{ weber } / 1 \text{ amp}
            .\end{align*}
        \item \textbf{Work}: Unsurprisingly, the energy stored in the magnetic field of an inductor is proportional to the inductance. It is also proportional to the square of the current through the inductor
            \begin{align*}
                W = \frac{1}{2}LI^{2}
            .\end{align*}
        \item \textbf{Simple inductor}: An inductor in its simplest form consists of a series of wire loops. These might be wound around an iron core, although a non-ferrous core might also be used. For a simple single layer inductor, the inductance is described by the following formula:
            \begin{align*}
                L = \mu \frac{AN^{2}}{l}
            .\end{align*}
            \begin{itemize}
                \item $L$ is the inductance in henries,
                \item $\mu$ is the permeability of the core material,
                \item $A$ is the cross sectional area of the coil,
                \item $N$ is the number of coils or turns,
                \item $l$ is the length of the coil.
            \end{itemize}
            \fig{.7}{./figures/103.png}
        \item \textbf{Permeability}: Inductor permeability is a measure of how easily a magnetic core material can concentrate and support magnetic flux when an electric current flows through the inductor's coil.
        \item \textbf{Why coils}: to concentrate and amplify the magnetic field, which maximizes the component's ability to store energy and oppose changes in electric current.
            \bigbreak \noindent 
            When electric current flows through a straight wire, it creates a weak, scattered circular magnetic field. By winding that same wire into a coil (solenoid), the magnetic fields from every individual loop add together (superpose) inside the center, creating a strong, dense, and uniform magnetic field.
        \item \textbf{Inductor Styles and Packaging}: In order to achieve high inductance, we would like a core with high permeability, permeability being a measure of how easy it is to establish magnetic flux in said material. Substances such as iron or ferrite have a much greater permeability than air and are used commonly for cores. They do have a disadvantage in that they will saturate sooner than an air core, and this can lead to distortion. 
            \bigbreak \noindent 
            Another approach is to pack as many turns as is possible within a given length. One way to do this is to minimize the thickness of the insulation around the wire. This can be achieved by using a thin enamel coating instead of the typical plastic insulation. 
            \bigbreak \noindent 
            A second method is to use a very fine wire. This leads to two problems, namely an undesirable increase in equivalent series resistance (known colloquially as coil resistance or $R_{\text{coil}}$), and limited current carrying capacity. All of these effects have to be balanced in order to achieve the best performance for a given application.
            \bigbreak \noindent 
            Commercial inductors range in value from a fraction of a nanohenry for small surface mount “chip” inductors up to several henries. Some devices exhibit large internal inductances even though they are not specifically used as inductors. One common example is a transformer. Another example is an electric guitar or bass pickup. Units such as this may be constructed of several thousand turns of very fine wire (typically AWG 41 to 44) and achieve inductances in excess of one henry.
            \bigbreak \noindent 
            \textbf{Note:} The wound wire must be insulated otherwise each loop will short to the loops next to it and we'd be left with a tube instead of a series of loops.
            \bigbreak \noindent 
            \fig{.7}{./figures/104.png}
            \bigbreak \noindent 
            A variety of inductors is shown above, all of which are of the through-hole type (surface mount inductors do not appear considerably different from their surface mount resistor and capacitor cousins).
            \bigbreak \noindent 
            The two units toward the left are molded inductors and use a standard color code, similar to the ones used for resistors and capacitors. The unit at the top (yellow) is a high current inductor that features low $R_{\text{coil}}$.
            \bigbreak \noindent 
            The three inductors in the center use obvious ferrite cores, two wound on straight cores and the third wound on a toroidal core. The unit to the right uses a high permeability material at the very top and is wrapped in a plastic sheath for protection. Variable inductors are also a possibility and can be made by using a ferrite core that can be slid within the coils, effectively changing the permeability of the core (part ferrite, part air).
        \item \textbf{Schematic symbol}: 
            \bigbreak \noindent 
            \fig{.8}{./figures/105.png}
            \bigbreak \noindent 
            The standard symbol is at the top. The variable inductor symbol is in the middle and is a two lead device, somewhat reminiscent of the symbol for a rheostat. At the bottom is the symbol for an inductor with an iron, ferrite, or similar high permeability core. In general, like resistors, single inductors are not polarized and cannot be inserted into a circuit backwards. There are, however, special applications where multiple coils can be wound on a common core, and for these, the polarity of their interconnection can make a difference.
        \item \textbf{Inductors in Series and in Parallel}: Suppose we take two identical inductors and place them in series. This effectively doubles both the length and the number of loops. 
            \bigbreak \noindent 
            Suppose we take two identical inductors and place them in series. This effectively doubles both the length and the number of loops. This is because $N^{2}$ goes up by a factor of four which is then halved be the increased length. Consequently, inductors in series add values just like resistors in series. By extension, inductors in parallel behave like resistors in parallel. The equivalent of parallel inductors can be found by using either the product-sum rule or by taking the reciprocal of the sum of their reciprocals. 
            \bigbreak \noindent 
            Consider
            \bigbreak \noindent 
            \fig{.8}{./figures/106.png}
            \bigbreak \noindent 
            We have
            \begin{align*}
                L = (12 || 6 + 5)mH = 9mH
            .\end{align*}
        \item \textbf{Current voltage relationship}: The fundamental current-voltage relationship of the inductor is the mirror image of that of the capacitor
            \begin{align*}
                v = L \frac{di}{dt}
            .\end{align*}
            This states that the voltage across the inductor is a function of how quickly the current is changing. If the current is not changing (i.e., in steady-state), then the voltage across the inductor is zero. In this case, the inductor behaves like a short, or more accurately, like its $R_{\text{coil}}$ value. In contrast, during a rapid initial current change, the inductor voltage can be large, and thus the inductor behaves like an open. In contrast, during a rapid initial current change, the inductor voltage can be large, and thus the inductor behaves like an open.
            \bigbreak \noindent 
            An inductor resists any change in its magnetic field. When current tries to rise rapidly from zero, the expanding magnetic field induces a voltage (back-EMF) that opposes the incoming source voltage.
            \bigbreak \noindent 
            Because the inductor drops nearly all of the supply voltage to force the initial current to start at zero, very little current flows right away. From the perspective of the rest of the circuit, it acts like an open circuit (infinite initial resistance/impedance).
            \bigbreak \noindent 
            When you first apply power to a circuit containing an inductor, the current tries to rush in. Because an inductor hates sudden changes in current, it fights back by instantly generating an opposing voltage (back-EMF) that is practically equal to the power supply's voltage.
            \bigbreak \noindent 
            An inductor resists changes in current because a changing current creates a changing magnetic field, which induces a fighting voltage (back-EMF) governed by Faraday's Law of Induction and Lenz's Law.  When current flows through an inductor’s coil, it builds a magnetic field. Changing that current means the magnetic field must grow or shrink. According to Faraday's Law, a changing magnetic field induces a voltage (electromotive force, or EMF) inside the coil. Lenz's Law states that this induced voltage will always try to oppose the change that caused it. If you try to increase the current, the inductor pushes back with a voltage that fights the incoming flow. If you try to decrease or stop the current, the collapsing magnetic field releases stored energy and pushes current forward to keep it from dropping.
            \bigbreak \noindent 
            If your battery is 10 Volts, the inductor instantly generates an opposing 10 Volts across itself. Because it matches the battery's push, it "absorbs" or drops the entire voltage right at its terminals. Because the inductor has generated an opposing voltage equal to the battery, the net voltage pushing current through the rest of the circuit is essentially zero (\(10\text{V} - 10\text{V} = 0\text{V}\)). With no net voltage to push it, the current is forced to start at exactly zero Amps at the very first moment (\(t = 0\)).
            \bigbreak \noindent 
            The inductor has "hogged" or absorbed all of the circuit's available electrical pressure. Because it uses up the full 10V just to fight the battery, there is no voltage left over to push current through the rest of the circuit.
            \bigbreak \noindent 
            Notice that
            \begin{align*}
                \frac{d}{dt} = \frac{v}{L}
            .\end{align*}
            Thus if an inductor is fed by a constant voltage source, the current will rise at a constant rate equal to $v/L$.
            \bigbreak \noindent 
            Consider
            \bigbreak \noindent 
            \fig{.6}{./figures/108.png}
            \bigbreak \noindent 
            We see a voltage source feeding a single inductor. If we were to plot the inductor's current over time, we would see something like
            \bigbreak \noindent 
            \fig{.6}{./figures/107.png}
            \bigbreak \noindent 
            As time progresses, the current through the inductor increases, flowing from top to bottom. With a theoretically perfect inductor and source, this would continue as long as the circuit was energized. In reality, this line would either begin to deflect horizontally as the source reached its limits, or the inductor would fail once its maximum current or power handling was reached. The slope of this line is dictated by the size of the applied voltage source and the inductance.
            \bigbreak \noindent 
            Consider
            \bigbreak \noindent 
            \fig{.6}{./figures/109.png}
            \bigbreak \noindent 
            We have
            \begin{align*}
                \frac{di}{dt} = \frac{v}{L} = \frac{10V}{50mH} = \frac{200A}{s}
            .\end{align*}
            As noted previously, if an inductor is driven by a fixed voltage source and ignoring $R_{\text{coil}}$, the current through it rises at the constant rate of $v/L$. This change in current through the inductor is not limitless. An instantaneous change requires that $di/dt$ is infinite, and thus, the voltage driving the inductor would also have to be infinite, which is a clear impossibility. Therefore we can state a particularly important characteristic of capacitors:
            \bigbreak \noindent 
            \begin{quote}
               "\textit{The current through an inductor cannot change instantaneously}" 
            \end{quote}
        \item \textbf{Initial and Steady-State Analysis of RL Circuits}: When analyzing resistor-inductor circuits, remember that current through an inductor cannot change instantaneously as this would require an infinite voltage source. When a circuit is first energized, the current through the inductor will still be zero, which is characteristic of opens. Once at steady-state, the current has leveled out and therefore the voltage across the inductor will approach zero, which is characteristic of shorts. Thus, we can state the general behavior of inductors at the beginning and ending of the charge cycle:
            \begin{itemize}
                \item For DC analysis, initially inductors appear as opens.
                \item At steady-state, inductors appear as shorts.
            \end{itemize}
            This is the opposite of what was seen with capacitors.
            \bigbreak \noindent 
            Consider
            \bigbreak \noindent 
            \fig{.8}{./figures/110.png}
            \bigbreak \noindent 
            Initially $L$ is open, leaving us with $R_1$ and $R_2$ in series with the source, E. At steady-state, $L$ shorts out, leaving $R_1$ in series with the parallel combination of $R_2$ and $R_3$. All practical inductors will exhibit some internal resistance, so it is often best to think of an inductor as an ideal inductance with a small resistance ($R_{\text{coil}}$) in series with it.
            \bigbreak \noindent 
            Now, consider 
            \bigbreak \noindent 
            \fig{.8}{./figures/111.png}
            \bigbreak \noindent 
            Initially, the inductor behaves as on open 
            \bigbreak \noindent 
            \fig{.7}{./figures/112.png}
            \bigbreak \noindent 
            Thus,
            \begin{align*}
                V_{2k} = I_{S}2k = \frac{25}{3k} \cdot 2k = 16.67V
            .\end{align*}
             For steady-state, we redraw using a short in place of the inductor
             \bigbreak \noindent 
             \fig{.8}{./figures/113.png}
             \bigbreak \noindent 
             Here,
             \begin{align*}
                 V_{2k} = I_{S}(6k || 2k) = \frac{25}{6k || 2k + 1k} \cdot 6k || 2k = 15V
             .\end{align*}
        \item \textbf{Initial and Steady-State Analysis of RLC Circuits}: When analyzing resistor-inductor-capacitor circuits, remember that capacitor voltage cannot change instantaneously, thus, initially, capacitors behave as a short circuit. Once the capacitor has been charged and is in a steady-state condition, it behaves like an open. This is opposite of the inductor. As we have seen, initially an inductor behaves like an open, but once steady-state is reached, it behaves like a short
            \bigbreak \noindent 
            \fig{.8}{./figures/114.png}
            \bigbreak \noindent 
            Initially $L$ is open and $C$ is a short, leaving us with $R_1$ and $R_2$ in series with the source, $E$. At steady-state, $L$ shorts out both $C$ and $R_2$, leaving all of $E$ to drop across $R_1$. For improved accuracy, replace the inductor with an ideal inductance in series with the corresponding $R_{\text{coil}}$ value. Similarly, practical capacitors can be thought of as an ideal capacitance in parallel with a very large (leakage) resistance.
            \bigbreak \noindent 
            Consider
            \bigbreak \noindent 
            \fig{.7}{./figures/115.png}
            \bigbreak \noindent 
            Assuming the initial current through the inductor is zero and the capacitor is uncharged, determine the current through the 2 k$\Omega$ resistor when power is applied and after the circuit has reached steady-state
            \bigbreak \noindent 
            \fig{.7}{./figures/116.png}
            \bigbreak \noindent 
            The shorted capacitor removes everything to its right from the circuit. All that's left is the source and the 2 k$\Omega$ resistor
            \bigbreak \noindent 
            We can find the current through the 2 k$\Omega$ resistor using Ohm's law
            \begin{align*}
                I_{2k} = \frac{14V}{2k\Omega} = 7mA
            .\end{align*}
            In steady-state, 
            \bigbreak \noindent 
            \fig{.8}{./figures/117.png}
            \bigbreak \noindent 
            We are left with a resistance of 2 k$\Omega$ in series with the parallel combination of 1 k$\Omega$ and 4 k$\Omega$, or 2.8 k$\Omega$ in total.
            \bigbreak \noindent 
            \begin{align*}
                I_{2k} = \frac{14V}{2.8k\Omega} = 5mA
            .\end{align*}
        \item \textbf{Transient Response of RL Circuits}: The transient response of RL circuits is nearly the mirror image of that for RC circuits.
            \bigbreak \noindent 
            \fig{.7}{./figures/118.png}
            \bigbreak \noindent 
            Again, the key to this analysis is to remember that inductor current cannot change instantaneously. When power is first applied, the circulating current must remain at zero. Therefore no voltage drop is produced across the resistor, and by KVL, the voltage across the inductor must equal the source, $E$. This establishes the initial rate of change of current via
            \begin{align*}
                \frac{di}{dt} = \frac{E}{L}
            \end{align*}
            and is represented by the dashed red line. As the current starts to increase, the voltage drop across the resistor begins to increase. This reduces the voltage available for the inductor, thus slowing the rate of change of current. This is depicted by the solid red curve on the graph.
            \bigbreak \noindent 
            Meanwhile, the solid blue curve represents the decreasing inductor voltage. Thus, in the RL circuit, the inductor's voltage curve echoes the RC circuit's current curve (or resistor voltage curve), and the RL current curve echoes the RC circuit's capacitor voltage curve.
            \bigbreak \noindent 
            \fig{.4}{./figures/119.png}
            \bigbreak \noindent 
            As noted before, the rate of current change versus time is equal to $v/L$, and therefore in this case, $E/L$. If the initial rate of change were to continue unabated, the maximum (steady-state) current, $E/R$, would be reached in $L/R$ seconds. Therefore the time constant for an RL circuit is
            \begin{align*}
                \tau = \frac{L}{R}
            .\end{align*}
            Once again, five constants will achieve steady-state. Following the prior work on capacitors, the relevant equations for the RL circuit can be shown to be
            \begin{align*}
                V_{L}(t) &= Ee^{-\frac{t}{\tau}}, \\
                V_{R}(t) &= E\left(1-e^{-\frac{t}{\tau}}\right), \\
                I(t) &= \frac{E}{R}\left(1-e^{-\frac{t}{\tau}}\right)
            .\end{align*}
            \bigbreak \noindent 
            Consider
            \bigbreak \noindent 
            \fig{.8}{./figures/120.png}
            \bigbreak \noindent 
            Assume the switch is closed at time $t = 0$. Determine the charging time constant, the amount of time after the switch is closed before the circuit reaches steady-state, and the inductor voltage and current at $t = 0$, $t = 2$ microseconds and $t = 1$ millisecond. Assume the inductor is initially uncharged.
            \bigbreak \noindent 
            We have
            \begin{align*}
                \tau = \frac{L}{R} = \frac{400\mu H}{150\Omega} = 2.667\mu s
            .\end{align*}
            Thus, steady-state is reached in 
            \begin{align*}
                5\tau = 2.667(5) = 13.33\mu s
            .\end{align*}
            So,
            \begin{align*}
                V_{L}(0) = 9V, \quad V_{L}(1) = 0V
            .\end{align*}
            Because the inductor is initially open, 
            \begin{align*}
                I_{L}(0) = 0A
            .\end{align*}
            At $I_{L}(1)$, the circuit is in steady-state and the inductor acts like a short. Thus, all of the 9 volt source drops across the 150 $\Omega$ resistor, for 60mA.
            \bigbreak \noindent 
            To find $V_{L}(2\mu s)$,
            \begin{align*}
                V_{L}(t) = Ee^{-\frac{t}{\tau}} \implies V_{L}(2\mu s) = 9Ve^{-\frac{2\mu s}{2.667\mu s}} = 4.251V
            .\end{align*}

    \end{itemize}

    \pagebreak 
    \unsect{Magnetic Circuits and Transformers}
    \begin{itemize}
        \item \textbf{Intro}: Here we shall take a closer look at the parameters of magnetic circuits such as permeability and flux density, and introduce new parameters including reluctance and magnetizing force. We shall also examine a key graphic that helps define the magnetic properties of materials, namely the BH curve.        
            \bigbreak \noindent 
            As we shall see, there are certain parallels between magnetic circuits and electrical circuits. For example, there is an “Ohm's law for magnetic circuits”, sometimes referred to as Hopkinson's law
            \bigbreak \noindent 
            There is also a version of Kirchhoff's voltage law dealing with the magnetic equivalents of voltage rises and drops.
            \bigbreak \noindent 
            Perhaps the two most common applications of magnetic circuits for the new student involve the use of transformers and relays.
            \bigbreak \noindent 
            Other common applications include dynamic loudspeakers and microphones, magnetic resonance imaging (i.e., medical MRI)
        \item \textbf{Electromagnetic Induction}: Perhaps the most important observation regarding magnetic systems is Faraday's law of electromagnetic induction. Briefly, it states:
            \begin{quote}
               "\textit{If a conductor is cut by changing magnetic lines of force, a voltage will be induced in the conductor}" 
            \end{quote}
            More specifically, 
            \begin{align*}
                e = -\frac{d\Phi}{dt}
            .\end{align*}
            Note that $e$ is the \textbf{induced voltage}, and $\frac{d\Phi}{dt}$ is the rate of change of magnetic flux with respect to time.
            \bigbreak \noindent 
            Therefore, the larger the flux and the more quickly it fluctuates, the greater the induced voltage. What is important here is a relative change with respect to the conductor and field. This can be accomplished in two basic ways: by a magnetic field that is itself fluctuating around a fixed conductor, or by a conductor moving through a magnetic field.
        \item \textbf{Transducer}: The concept of electromagnetic induction is a key element of a variety of transducers. Simply put,
            \begin{quote}
                "\textit{A transducer is a device that transforms energy from one type to another.}" 
            \end{quote}
            An electric motor transforms electrical energy into mechanical energy and a generator does the precise opposite, transforming mechanical energy into electrical energy. More commonly, the word transducer is associated with devices such as loudspeakers and microphones. A loudspeaker transforms its electrical input into an acoustic pressure wave (sound) whereas the microphone performs the complementary function of transforming an acoustic pressure wave into an electrical signal. While there are different ways of constructing these devices, the most common designs are the dynamic loudspeaker and dynamic microphone, both of which rely on the principle of electromagnetic induction.
        \item \textbf{Dynamic Loudspeakers and Microphones}: Regardless of their output capabilities and frequency range, all dynamic loudspeakers share a common set of elements. Corresponding elements can be found in dynamic microphones. Indeed, the devices are so similar that in some applications, small loudspeakers might pull double duty and switch over to a microphone mode of operation.
            \bigbreak \noindent 
            \fig{.4}{./figures/121.png}
            \bigbreak \noindent 
            The heart of the system is a voice coil (H) which is a coil of tightly wound magnet wire. This sits inside of a magnetic field that is created by a powerful permanent magnet (E). The voice coil is connected to a diaphragm (F) that is is usually made of paper or plastic. This assembly is connected to the frame by the spider (D) at the top of the voice coil, and a suspension at the end of the diaphragm (B). The voice coil and diaphragm can move back and forth relative to the frame, like a piston.
            \bigbreak \noindent 
            To understand the operation, recall that passing a current through a coil of wire creates a magnetic field around the coil
            \bigbreak \noindent 
            This magnetic field is the same as is created by a normal magnet. In this case the north pole is on the left side (i.e., lines of force exiting). Consequently, this coil will interact with a permanent magnet just like any other magnet would. In the case of the loudspeaker, the voice coil is fed a current from the amplifier that echoes the music or voice signal. This current creates a magnetic field around the voice coil which interacts with the field of the permanent magnet. As the direction of the current changes, the poles of the voice coil's field reverse. Thus, sometimes the voice coil is pushed outward and sometimes it is pulled inward. Further, the larger the current, the greater the field it creates, and the greater the push or pull.
            \bigbreak \noindent 
            \fig{.6}{./figures/122.png}
            \bigbreak \noindent 
            This results in a back-and-forth motion of the diaphragm that echoes the shape of the waveform being fed from the amplifier. As the diaphragm pushes against the air, it sets up a pressure wave, and we have sound.
            \bigbreak \noindent 
            A dynamic microphone runs the sequence in reverse. First, the diaphragm will move back and forth in accordance to an applied sound wave. This causes the voice coil to move back and forth within the field of the permanent magnet. We now have a coil of wire being cut by magnetic lines of force, and by Faraday's law of induction, this means that a voltage will be induced in the coil. As long as the diaphragm can respond to the subtle changes in the acoustic pressure wave, then the resulting induced voltage should be of high fidelity. Clearly, to do this with great accuracy, the diaphragm/voice coil/suspension assembly will need to very light and nimble, meaning that the components of the microphone will tend to be much smaller than those of a loudspeaker.
        \item \textbf{Magnetic Circuits}: Magnetic circuits include applications such as transformers and relays. A very simple magnetic circuit is shown below
            \bigbreak \noindent 
            \fig{.8}{./figures/123.png}
            \bigbreak \noindent 
            First, it consists of a magnetic core. The core may be comprised of a single material such as sheet steel but can also use multiple sections and air gap(s). Around the core is at least one set of turns of wire, i.e., a coil formed around the core. Multiple sets of turns are used for transformers
            \bigbreak \noindent 
            As we have seen already, passing current through the windings generates a magnetic flux, $\Phi$, in the core. As this flux is constrained within the cross sectional area of the core, $A$, we can derive a flux density, $B$.
            \begin{align*}
                B = \frac{\Phi}{A}
            .\end{align*}
            \begin{itemize}
                \item $B$ is the magnetic flux density in teslas,
                \item $\Phi$ is the magnetic flux in webers,
                \item $A$ is the area in square meters.
            \end{itemize}
            Recall that one tesla is defined as one weber per square meter. An alternate unit that is sometimes used is \textbf{gauss}.
            \begin{align*}
                1 \text{ tesla } = 10^{3} \text{ gauss }
            .\end{align*}
            Consider an example. Suppose a magnetic flux of $6\times 10^{-5}$ webers exists in a core whose cross section has dimensions of 1 centimeter by 2 centimeters. Determine the flux density in teslas.
            \begin{align*}
                B = \frac{\Phi}{A} = \frac{6\times 10^{-5}}{0.01m \times 0.02m} = 0.3 T
            .\end{align*}
        \item \textbf{Ohm's Law for Magnetic Circuits (Hopkinson's or Rowland's Law)}: There is a common parallel drawn between magnetic circuits and electrical circuits, namely Hopkinson's law.
            \begin{align*}
                F = \Phi R
            .\end{align*}
            \begin{itemize}
                \item $F$ is the magnetomotive force (or MMF) in amp-turns,
                \item $\Phi$ is the magnetic flux in webers,
                \item $R$ is the reluctance of the material in amp-turns/weber
            \end{itemize}
            The magnetomotive force compares to a source voltage or electromotive force (EMF), the magnetic flux is likened to the flow of current, and the reluctance stands in for resistance (i.e., on the one hand we have a material that resists the flow of current, and on the other, a material in which there is “reluctance” to establish magnetic flux). Further, magnetomotive force is the product of the current flowing through a coil and the number of loops or turns in the coil:
            \begin{align*}
                F = NI
            .\end{align*}
            \begin{itemize}
                \item $F$ is the magnetomotive force in amp-turns,
                \item $N$ is the number of loops or turns in the coil,
                \item $I$ is the current in the coil in amps.
            \end{itemize}
        \item \textbf{amp-turns}: The unit of magnetomotive force (MMF), which is the magnetic "pressure" or driving force that produces a magnetic field in an electromagnetic coil
            \bigbreak \noindent 
            One ampere-turn is equal to the magnetomotive force produced when a direct current of one ampere flows through a single-turn loop of wire
        \item \textbf{Reluctance}: The equation for reluctance has a nice parallel with the equation for resistance
            \bigbreak \noindent 
            \begin{align*}
                R = \frac{l}{\mu A}
            .\end{align*}
            \begin{itemize}
                \item $R$ is the reluctance in amp-turns/weber,
                \item $l$ is the length of the material in meters,
                \item $A$ is the cross sectional area of the material in square meters,
                \item $\mu$ is the permeability of the material in henries/meter.
            \end{itemize}
        \item \textbf{Magnetizing force}: Given the characteristics of the coil and the path length of the magnetic circuit, the magnetic flux gives rise to a magnetizing force, $H$. 
            \begin{align*}
                H = \frac{NI}{l}
            .\end{align*}
            \begin{itemize}
                \item $H$ is the magnetizing force in amp-turns/meter,
                \item $N$ is the number of turns or loops in the coil,
                \item $I$ is the coil current in amps,
                \item $l$ is the length of the magnetic path in meters.
            \end{itemize}
        \item \textbf{Relation}: Ferromagnetic materials (i.e., materials that have high permeability, such as steel) produce a low reluctance. The practical problem here is that $\mu$, unlike the resistivity $\rho$ of resistors, is not a constant for these sorts of materials. It can vary considerably.
            \bigbreak \noindent 
            \fig{.7}{./figures/124.png}
            \bigbreak \noindent 
            As a result, it is impractical to find reluctance in same manner that we find resistance. All is not lost, though.
            \bigbreak \noindent 
            The flux density and corresponding magnetizing force for any given material are related by the following equation
            \begin{align*}
                B = \mu H
            .\end{align*}
            \begin{itemize}
                \item $B$ is the flux density in teslas,
                \item $\mu$ is the permeability of the material in henries/meter,
                \item $H$ is the magnetizing force in amp-turns/meter.
            \end{itemize}
            Once again, the tricky bit here is the permeability of the core material. For air, we can use the permeability of free space, $\mu_{0}$
            \begin{align*}
                \mu_{0} = 4\pi \times 10^{-7} \frac{H}{m} \approx 1.257 \times 10^{-6} \frac{H}{m}
            .\end{align*}
        \item \textbf{BH curve intro}: For other materials, such as sheet steel or cast steel, we shall take another path; namely an empirically derived curve that plots flux density $B$ against magnetizing force $H$
            \bigbreak \noindent 
            \fig{.7}{./figures/125.png}
            \bigbreak \noindent 
            Clearly, this curve is not a nice straight line, or even an obvious, predictable function. The immediately apparent attributes are the initial steep rise which is followed by a flattening of the curve. This flattening corresponds to saturation of the magnetic material. In contrast, a plot for air would reveal a straight line with a very shallow slope. As we shall see, being able to achieve a high flux density for a given magnetizing force will result in an effective and efficient magnetic circuit. So while air has the positive attribute of not saturating, the resulting flux density is low, usually leading to lower performance.
        \item \textbf{The BH curve}:  The process of generating a BH curve is as follows. First, we create a core of the material to be investigated. A coil of wire is then wrapped around this core
            \bigbreak \noindent 
            \fig{.7}{./figures/126.png}
            \bigbreak \noindent 
            Here we have a basic toroid with a coil of $N$ turns. We begin with the system at rest and not energized. A small current, I, is applied to the coil. This produces a magnetizing force via
            \begin{align*}
                H = \frac{NI}{l}
            .\end{align*}
            There will be a corresponding flux density, as per
            \begin{align*}
                B = \mu H
            .\end{align*}
            The current is then increased. This produces an increase in magnetizing force and a corresponding change in flux density. The current is increased further, until the curve flattens, indicating that saturation has been reached.
            \bigbreak \noindent 
            \fig{.4}{./figures/126.png}
            \bigbreak \noindent
            Starting at point $a$ with point $b$ indicating saturation. The current is then reduced. This causes a reduction in magnetizing force, but while the flux density decreases, it does not trace back perfectly along the original trajectory. Instead, the curve is displaced above the original.
            \bigbreak \noindent 
            Eventually, the current will be reduced to zero. This corresponds to point $c$. At this point, even though there is no current in the coil, there is still flux in the core. The resulting flux density is referred to as \textbf{retentivity} and is a measure of residual magnetism. This phenomenon is what makes it possible to magnetize materials. 
            \bigbreak \noindent 
            If we now reverse the direction of the coil current and begin to increase its magnitude, the flux density will continue to drop. At point $d$, it will have reached zero. As we have effectively coerced the flux back to zero, we call the magnetizing force required to do this the \textbf{coercivity} or \textbf{coercive force}.
            \bigbreak \noindent 
            As the current magnitude increases, the flux density also increases but with opposite sign. Eventually, saturation is reached again at point $e$. Once again, if the current's magnitude is decreased, the magnitude of the flux density will also decrease but will not perfectly trace back along the original trajectory. This time it will follow a lower path. When the current is reduced to zero at point $f$, a mirror retentivity is evident. Further positive increases in current show a mirror coercivity at point $g$. Finally, as current is increased to maximum, we again reach saturation at point $b$.
            \bigbreak \noindent 
            If the current is cycled in this manner again, the process is repeated, and the outer path indicated by the arrows is taken again. Thus, a specific value of magnetizing force can give rise to different values of flux density: it depends on the recent history of the material. This effect is known as hysteresis and is found in other areas as well.
            \bigbreak \noindent 
            Effectively, published BH curves follow the middle of the hysteresis loop. An example of BH curves for three different core materials is shown
            \bigbreak \noindent 
            Curve A is for sheet steel (which is common in transformers), curve B is for cast steel, and curve C is for cast iron. We shall make use of these in upcoming examples. Curves for other materials are also available.
            \bigbreak \noindent 
            \fig{.4}{./figures/128.png}
            \bigbreak \noindent 
            As an example, assume the toroid we saw earlier is made of cast steel, has a 500 turn coil, a cross section of 2 cm by 2 cm, and an average path length of 50 cm. Determine the flux in webers if a current of 0.3 amps feeds the coil.
            \bigbreak \noindent 
            So, 
            \begin{align*}
                H = \frac{NI}{l} = \frac{500(0.3)}{0.5}  = 300 \frac{At}{m}
            .\end{align*}
            Using the curve for cast steel, a decent for the flux density is 0.52 teslas. This is the corresponding flux density. In order to find the flux, we need to find the area of the core.
            \begin{align*}
                A = 0.02^{2} = 0.0004m^{2}
            .\end{align*}
            Thus,
            \begin{align*}
                \Phi = BA = 0.52T \times 0.0004m^{2} = 2.08 \times 10^{-4}Wb
            .\end{align*}
        \item \textbf{Summary}: 
            \begin{itemize}
                \item \textbf{Electromagnetic induction}: Electromagnetic induction is the process of generating a voltage and electric current in a conductor by placing it in a changing magnetic field or moving it relative to one
                    \begin{align*}
                        e = -\frac{d\Phi}{dt}
                    .\end{align*}
                    $e$ is the \textbf{induced voltage}.
                \item \textbf{Magnetic flux ($\Phi$)}: Magnetic flux is the total measure of a magnetic field passing through a given surface area
                    \bigbreak \noindent 
                    Measured in \textbf{Weber (Wb)}
                \item \textbf{Magnetic flux density}:  Magnetic flux density is a vector quantity that measures the strength and direction of a magnetic field based on the number of magnetic field lines passing through a given area
                    \begin{align*}
                        B = \frac{\Phi}{A} = \mu H
                    .\end{align*}
                    Measured in \textbf{Telsa (T)} or \textbf{Gauss (G)}.
                \item \textbf{Magnetomotive force (amp-turns)}: Magnetomotive force (MMF) is the magnetic "pressure" or potential that drives magnetic flux through a magnetic circuit
                    \bigbreak \noindent 
                    You can think of MMF as the magnetic equivalent of voltage (electromotive force, or EMF) in an electrical circuit. Just as voltage pushes electric current through a wire, MMF pushes magnetic field lines (flux) through a core or material
                    \begin{align*}
                        F = NI
                    \end{align*}
                    where $N$ is the number of turns (loops) in a wire coil, and $I$ is the current through the wire. Magnetomotive force is measured in \textbf{amp-turn (At)}.
                \item \textbf{Magnetic Reluctance}:  The opposition a material offers to magnetic flux, measured in ampere-turns per weber
                    \begin{align*}
                        R = \frac{l}{\mu A}
                    ,\end{align*}
                    where $l$ is the path length, $A$ is the cross-sectional area, and $\mu$ is the permeability of the material
                    \bigbreak \noindent 
                    Unlike electrical resistance, magnetic reluctance does not dissipate power or energy as heat; it merely stores energy in the magnetic field
                \item \textbf{Hopkinson's Law}: Hopkinson's law is the magnetic circuit counterpart to Ohm's law, stating that the magnetomotive force is equal to the magnetic flux multiplied by the magnetic reluctance
                    \begin{align*}
                        F = \Phi R
                    .\end{align*}
                \item \textbf{Magnetizing force}: The magnetizing force ($H$), also known as magnetic field intensity, is the external magnetic field strength applied to a material to induce magnetic flux or align internal magnetic domains
                    \begin{align*}
                       H = \frac{NI}{l} 
                    \end{align*}
                    Where $I$ is current, $N$ is the number of coil turns, and $l$ is the length of the magnetic circuit. 
                    \bigbreak \noindent 
                    Measured in amp-turns per meter ($\frac{At}{m}$)
                \item \textbf{Magnetizing force vs magnetomotive force}: Magnetomotive force (MMF) is the total "pressure" or driving potential that creates magnetic flux in a magnetic circuit, while magnetizing force (\(H\)) is the magnetomotive force distributed per unit length of the magnetic path (also called magnetic field strength).
            \end{itemize}
        \item \textbf{KVL for Magnetic Circuits}: We have
            \begin{align*}
                F = Hl = NI
            .\end{align*}




    \end{itemize}









    
\end{document}
